\section{The arithmetic count}\label{sec:A}
We write
\[
 F(s,t,z)=s^2-3t^2-18z^2,\quad
 \varepsilon_0=2+\sqrt3,\quad L_0=\log\varepsilon_0,
 \quad b_0=\log\frac3{2\sqrt2}.
\]
The letter $t$ in this section is an integer coordinate, unrelated to the
elliptic-surface parameter in Part C. All areas on the hyperbolic plane use
curvature $-1$ and the measure $dx\,dy/y^2$.

\subsection{The coordinates, ideals, and all boundary terms}
\paragraph{Statement A.1.1 (Proved).}
For every real $X\geq1$, the count in the prompt equals
\[
 \#\{(s,t,z)\in\mathbb Z^3:F(s,t,z)=-2,\ 1\leq z\leq X,
 \ t\geq0,\ t+4z<s\leq17z/4\}.
\]
The comparison map from the original variables is the integral isometry
\[
 \begin{pmatrix}s\\t\\z\end{pmatrix}=
 \begin{pmatrix}17&0&72\\0&1&0\\4&0&17\end{pmatrix}
 \begin{pmatrix}a\\q\\n\end{pmatrix},\qquad
 (a,q,n)=(17s-72z,t,17z-4s).
\]
It is a bijection of the indicated integer sets, with no parity restriction
lost. In the original $(x,y,z)$ coordinates,
\[
 s=x+y+4z,\quad t=y-x,\quad
 x=(s-4z-t)/2,\quad y=(s-4z+t)/2,
\]
and $2\{4(xy+xz+yz)-x^2-y^2-z^2\}=F(s,t,z)$.
\begin{proof}
The matrix has determinant $1$ and preserves $\operatorname{diag}(1,-3,-18)$.
The upper inequality is precisely $n\geq0$, and the strict lower inequality
is precisely $a+4n>q$. Conversely these inequalities force $a>0$:
if $a\leq0$ and $n=0$ the strict inequality is impossible; if $n\geq1$, it
implies $-4n<a\leq0$, whereas the equation gives
$a^2=3q^2+18n^2-2\geq16n^2$, a contradiction. The transformed height is
exactly $z=4a+17n$. Reduction of $F=-2$ modulo $2$ gives $s\equiv t\pmod2$,
so the displayed inverse to $(x,y,z)$ is integral. Direct expansion proves
its polynomial identity, and proves every asserted equivalence.
\end{proof}

\paragraph{Statement A.1.2 (Proved).}
Every nonzero integral ideal of $\mathbb Z[\sqrt3]$, of absolute ideal
norm $m$, has a unique generator $\alpha=s+t\sqrt3$ satisfying
\[
 \sqrt m\leq\alpha<\varepsilon_0\sqrt m.
\]
Its field norm may be $+m$ or $-m$. For example $(1+\sqrt3)$ has ideal
norm $2$ and field norm $-2$; a generator of field norm $+2$ is impossible
modulo $3$. Thus the interval assertion in Lead 2 is correct as written,
but it must not be silently strengthened to positive field norm for every
$m$. For the norms actually used here,
\[
 m_z=18z^2-2\quad(z\geq1),
\]
every ideal of norm $m_z$ has positive field-norm generator, so the unique
interval generator satisfies
\[
 N_{\mathbb Q(\sqrt3)/\mathbb Q}(\alpha)=m_z,\qquad s>0,\quad t\geq0.
\]
The number of these ideals is
\[
 \rho(m_z)=\sum_{d\mid m_z}\chi_{12}(d)
          =\rho(3z-1)\rho(3z+1).
\]
\begin{proof}
The full ring of integers is $\mathbb Z[\sqrt3]$. It is norm Euclidean:
rounding the two coefficients of an element of $\mathbb Q(\sqrt3)$ gives
an error $u+v\sqrt3$ with $|u|,|v|\leq1/2$ and
$|u^2-3v^2|\leq3/4<1$. Thus all ideals are principal. The units are
$\{\pm\varepsilon_0^k:k\in\mathbb Z\}$: a unit of norm $-1$ is impossible
modulo $3$, and the least unit greater than $1$ with positive conjugate is
$2+\sqrt3$. The latter assertion also follows directly since
$1<\alpha<\varepsilon_0$ would give an integral coefficient
$0< (\alpha-\alpha^{-1})/(2\sqrt3)<1$.

For any $m$, choose a generator, change its real sign to positive, and
multiply by a unique power of $\varepsilon_0$ to put it in the stated
half-open interval. This proves the general assertion without any
assumption on the sign of its field norm.
A generator of an ideal of norm $m_z$ has field norm $\pm m_z$.
Since $m_z\equiv1\pmod3$, the negative sign is impossible modulo $3$.
Choose its positive sign in the real embedding and multiply by a unique
power of $\varepsilon_0$ to put it in the stated half-open interval.
Then $s>0$ and $t\geq0$. The standard ideal factorization is particularly
transparent here: a split prime contributes $e+1$ ideals of norm $p^e$,
an inert prime contributes one if $e$ is even and none otherwise, and a
ramified prime contributes one. These are exactly the Euler factors of
$1*\chi_{12}$. The two linear factors $3z-1$ and $3z+1$ have greatest common
divisor dividing $2$; the extra factor $2$ and their common factors at $2$
make no difference because $2$ is ramified and $\rho(2^e)=1$.
\end{proof}

\paragraph{Statement A.1.3 (Proved).}
Put $\theta(\mathfrak a)=L_0^{-1}\log(\alpha/\sqrt{m_z})$ for the generator
in Statement A.1.2 and $\lambda_0=b_0/L_0$. For $X\geq1$,
\begin{equation}\label{A:arc-exact}
 R(X)=\sum_{1\leq z\leq X}
 \#\{\mathfrak a:N\mathfrak a=m_z,\ 0\leq\theta(\mathfrak a)\leq\lambda_0\}
 +R_0(X)-1,
\end{equation}
where
\[
 R_0(X)=\#\{k\geq1:4a_k\leq X\},\qquad
 a_k+q_k\sqrt3=(1+\sqrt3)\varepsilon_0^k.
\]
In particular $R_0(X)=O(\log(2X))$. The only points with $z\geq1$,
$t\geq0$, $s>0$, $s\leq17z/4$, and $F=-2$ which fail $s-4z>t$ are
$(4,0,1)$ and $(17,1,4)$.
\begin{proof}
Since
\[
 \cosh b_0=\frac{17}{12\sqrt2},
\]
the fixed-arc condition is exactly
\begin{equation}\label{A:arc-integer}
 288s^2\leq289(18z^2-2).
\end{equation}
It implies $4s<17z$. Conversely let $r=17z-4s\geq1$. Then
$r<17z$ and
\[
 18r(34z-r)\geq18(34z-1)\geq594>578,
\]
which is precisely the strict version of \eqref{A:arc-integer} after
substitution. Consequently the only points between the fixed arc and
the moving upper wall are the points $4s=17z$.

If $s-4z\leq t$, the equation implies
$(6z-s)^2\leq3z^2+1$. The upper wall gives
$(6z-s)^2\geq49z^2/16$, so $z\leq4$. For $z=1,2,3,4$ the interval
$\lceil\sqrt{18z^2-2}\rceil\leq s\leq\lfloor17z/4\rfloor$
gives respectively $(s,t)=(4,0)$, no point, no point, and $(17,1)$.
The first lies in the fixed arc; the second is on the moving wall.

On that wall write $z=4a$, $s=17a$. Then $a^2-3t^2=-2$.
Every solution with $a>0,t\geq0$ is
$(a,t)=(a_k,q_k)$ for a unique $k\geq0$. Indeed, reduction by the unit to
$\sqrt2\leq a+t\sqrt3<\varepsilon_0\sqrt2$ gives
$a\geq0$, $2a<3t$, and hence $3t^2<8$; its only possibility is $(1,1)$.
The strict wall excludes $k=0$ and includes all $k\geq1$.
Finally $a_k$ grows geometrically, proving the logarithmic estimate and
\eqref{A:arc-exact}, including both exceptional points.
\end{proof}

\paragraph{Statement A.1.4 (Proved by exact exhaustive computation).}
The independently computed counts are
\[
 \begin{array}{c|r}
 X&R(X)\\\hline
 10^6&35\,381\\
 10^7&353\,598\\
 10^8&3\,536\,002\\
 10^9&35\,363\,970
 \end{array}
\]
The first three confirm Lead 5. The fourth extends the census by a factor
of ten. The exact algorithm, its exhaustiveness proof, and its numerical
discrepancy analysis are given below in Statement A.5.1.

\subsection{The full integral group and the two closed geodesics}
Let $B_F(x,y)=x^{\mathsf T}\operatorname{diag}(1,-3,-18)y$ and
\[
 \mathbb H_F=\{x:F(x)=1,\ x_1>0\},\qquad
 O^+(F,\mathbb Z)=\{g\in O(F,\mathbb Z):g\mathbb H_F=\mathbb H_F\}.
\]
The orientation-preserving Fuchsian group is
$\Gamma^+=O^+(F,\mathbb Z)\cap SO(F,\mathbb R)$.
Thus the literal $SO(F,\mathbb Z)$ in Lead 3 must be restricted to its
index-two subgroup preserving the chosen component before it acts on
$\mathbb H_F$.

\paragraph{Statement A.2.1 (Proved).}
For the five primitive vectors
\[
 r_0=(0,-1,0),\ r_1=(0,0,-1),\ r_2=(1,1,0),\
 r_3=(9,3,2),\ r_4=(4,0,1),
\]
let $S_jx=x-2B_F(x,r_j)r_j/F(r_j)$. Then
$O^+(F,\mathbb Z)=\langle S_0,S_1,S_2,S_3,S_4\rangle$.
A fundamental domain is the compact right-angled pentagon
\[
 \mathcal P=\{(s,t,z)\in\mathbb H_F:
 t,z\geq0,\ s\geq3t,\ s\geq t+4z,\ 2s\geq9z\}.
\]
Its area is $\pi/2$. The group $\Gamma^+$ has area $\pi$, is generated by
$S_0S_j$ ($1\leq j\leq4$), and has signature $(0;2,2,2,2,2)$.
\begin{proof}
The five reflection matrices, in their displayed order, are
\[
 \begin{gathered}
 S_0=\operatorname{diag}(1,-1,1),\quad
 S_1=\operatorname{diag}(1,1,-1),\quad
 S_2=\begin{pmatrix}2&-3&0\\1&-2&0\\0&0&1\end{pmatrix},\\
 S_3=\begin{pmatrix}10&-9&-36\\3&-2&-12\\2&-2&-7\end{pmatrix},\qquad
 S_4=\begin{pmatrix}17&0&-72\\0&1&0\\4&0&-17\end{pmatrix}.
 \end{gathered}
\]
They are integral isometries, have determinant $-1$, and preserve
$\mathbb H_F$. In projective coordinates $(u,v)=(t/s,z/s)$ the pentagon
has vertices
\[
 (0,0),\ (1/3,0),\ (1/3,1/6),\ (1/9,2/9),\ (0,2/9).
\]
Each lies strictly inside $3u^2+18v^2<1$. Consecutive normals are
orthogonal; hence all five angles are $\pi/2$. Reflection across the sides
produces a tiling: glue copies along the paired sides, with four copies
around every vertex. The relations $S_j^2=1$ and
$(S_jS_{j+1})^2=1$ (indices modulo $5$) remove the edge and vertex loops,
so the resulting simply connected development is a hyperbolic surface.
It is complete: a bounded path cannot cross infinitely many nonadjacent
sides, since the pentagon is compact and these sides are separated by a
positive distance; only four chambers meet at a vertex. Its developing
map to the complete simply connected plane is consequently an isometry.
Thus the interiors of the pentagons are disjoint. The area formula gives
$3\pi-5\pi/2=\pi/2$.

We also prove that no integral isometries have been omitted. Let
$x=(a,q,n)\in\mathbb Z^3$ have $F(x)=1$, $a>0$. Reflections $S_0,S_1$
make $q,n\geq0$. If one of the other three defining inequalities fails,
reflection in that side strictly decreases the positive integer $a$.
It remains positive because the reflection preserves $\mathbb H_F$.
The process terminates in $\mathcal P$. There
\[
 a^2=(1-3u^2-18v^2)^{-1}\leq27/2;
\]
the maximum follows by convexity of $3u^2+18v^2$ on the projective polygon
and evaluation at its vertices. For integral $a\leq3$, the only norm-one
point in the pentagon is $(1,0,0)$: $a=2$ would give $(q,n)=(1,0)$,
which violates $a\geq3q$, and $a=3$ is impossible modulo $3$.

Apply this descent to $g(1,0,0)$ for any $g\in O^+(F,\mathbb Z)$.
After multiplication by a word in the $S_j$, it fixes $(1,0,0)$.
An integral isometry of the remaining positive binary form
$3q^2+18n^2$ only changes the two signs: its first column of norm $3$
is $(\pm1,0)$, and its second column is then $(0,\pm1)$.
These sign changes are $S_0,S_1$. This proves the asserted equality.
The orientation subgroup has index two, and doubling the pentagon gives
its stated orbifold signature and area.
\end{proof}

\paragraph{Statement A.2.2 (Proved).}
The integral vectors of norm $-2$ form exactly two $\Gamma^+$-orbits,
represented by $v_1=r_2=(1,1,0)$ and $v_2=r_4=(4,0,1)$.
They are distinguished by
\[
 \operatorname{div}(s,t,z)=\gcd(s,3t,18z),
\]
which equals $1$ and $2$, respectively. The stabilizers of the oriented
normals are primitive infinite cyclic groups. The associated primitive
closed-geodesic lengths are $2L_0$ and $L_0$. The oriented stabilizer of
$e_z=(0,0,1)$ is generated by
\[
 T=\begin{pmatrix}2&3&0\\1&2&0\\0&0&1\end{pmatrix}
\]
and its primitive length is $L_0$.
\begin{proof}
Move a point of the geodesic $v^\perp\cap\mathbb H_F$ into $\mathcal P$.
If $p$ is the resulting point and $B_F(p,v)=0$, the negative-definite
Cauchy--Schwarz inequality on $p^\perp$ gives
\[
 |v_1|^2=|B_F(e_s,v)|^2\leq2(p_1^2-1)\leq25.
\]
For $|s|\leq5$, the norm equation has only the possibilities
\[(s,t,z)=(\pm1,\pm1,0),(\pm4,0,\pm1),(\pm5,\pm3,0).\]
Checking the linear function $s-3tu-18zv$ at the five projective vertices
shows that only $\pm r_2,\pm r_4$ have geodesics meeting the pentagon.
These are exactly its two norm-$-2$ sides. A side reflection negates its
normal; an adjacent side reflection fixes it. Their product is in
$\Gamma^+$, so signs do not split the orbits. Moreover a determinant-$-1$
word mapping a representative to a vector can be multiplied by the
adjacent reflection fixing the representative. Thus passage from the
full reflection group to $\Gamma^+$ does not split either orbit.
The invariant divisibility separates them.

On any side, its two adjacent reflections act on its line by reflections
in the endpoints of a fundamental segment. The line stabilizer is the
resulting infinite dihedral group; its translation subgroup is generated
by their product. This follows directly by moving a point into that
segment in the reflection tiling; an additional stabilizer would identify
two interiors of fundamental segments. For $r_1,r_2,r_4$ the products are,
respectively, $S_0S_2$, $S_1S_3$, $S_3S_0$. Their traces are $5,15,5$.
For a hyperbolic isometry in the three-dimensional Lorentz representation,
$\operatorname{tr}g=1+2\cosh\ell(g)$. Since
$\cosh L_0=2$ and $\cosh(2L_0)=7$, the lengths follow. The source
stabilizer also follows immediately from the Pell equation
$u^2-3v^2=1$ for its upper left $2\times2$ block.
\end{proof}

\paragraph{Statement A.2.3 (Proved).}
For $v=(s,t,z)$ with $F(v)=-2$, $s>0$, $z\geq1$, let
$g_v=v^\perp\cap\mathbb H_F$ and $g_0=e_z^\perp\cap\mathbb H_F$.
They are disjoint and
\[
 \cosh d(g_v,g_0)=3z.
\]
Writing $m_z=s^2-3t^2$, the foot of their common perpendicular on $g_0$
has arclength coordinate
\[
 u=\log\frac{s+\sqrt3t}{\sqrt{m_z}},\qquad
 p(u)=(\cosh u,\sinh u/\sqrt3,0).
\]
In particular $\theta=u/L_0$ modulo $1$, and the arc in
\eqref{A:arc-exact} has length $b_0$ in the source circle of length $L_0$.
If $H_0=\langle T\rangle$ and $H_j=\operatorname{Stab}_{\Gamma^+}(v_j)$,
then, with cyclic stabilizers and both normal orientations retained,
\begin{equation}\label{A:double-coset-exact}
 \sum_{j=1}^2
 \#\{H_0\gamma H_j:\ 3|z(\gamma v_j)|\leq3X\}
 =4\sum_{1\leq z\leq X}\rho(m_z)+2.
\end{equation}
\begin{proof}
The unit spacelike normals have absolute inner product
$|B_F(v,e_z)|/\sqrt{(-2)(-18)}=3z>1$; the hyperboloid formula for
ultraparallel geodesics therefore gives the distance. To verify this
formula directly, restrict to their common two-dimensional Lorentz
plane, in which unit normals have coordinates $(0,1)$ and
$(\sinh d,\cosh d)$.

The minimum along $g_0$ occurs where
$s\sinh u-\sqrt3t\cosh u=0$. Solving gives the stated logarithm;
$T$ translates $u$ by $L_0$. The map
$\Gamma^+/H_j\to\Gamma^+v_j$, $\gamma H_j\mapsto\gamma v_j$,
is bijective by its stabilizer definition. Quotienting on the left by
$H_0$ gives precisely unit orbits of the norm equation in $(s,t)$.
For each $z\geq1$, there are $\rho(m_z)$ with $s>0$; the independent signs
of $z$ and $s$ give four such sets. At $z=0$, the norm-$-2$ equation has
two $H_0$-orbits, according to the sign of $s+\sqrt3t$. These account for
the last $2$ in \eqref{A:double-coset-exact}. This also tracks the
orbifold half-turns: they are not silently included in the cyclic
stabilizers or divided out a second time.
\end{proof}

\subsection{Comparison with the quaternion orders}
\paragraph{Statement A.2.4 (Proved).}
Let $B=(3,-1)_{\mathbb Q}$ with $i^2=3$, $j^2=-1$, $ij=k=-ji$.
Then $B$ has discriminant $6$. The orders
\[
 \mathcal O_{12}=\mathbb Z[1,i,j,k],\qquad
 \mathcal O_6=\mathbb Z[1,i,j,(1+i+j+k)/2],\qquad
 \mathcal O_{108}=\mathbb Z[1,i,3j,3k]
\]
have the indicated reduced discriminants, and $\mathcal O_6$ is maximal.
Under the pure-coordinate map
\[
 V(a,q,n)=\frac{3q+a}{2}i+\frac{3(q+a)}2j+3nk,
 \qquad \operatorname{nrd}V=\frac32F(a,q,n),
\]
the spinor kernel of $\Gamma^+$ is
$\operatorname{Ad}(\mathcal O_{12}^{1})$, of index $2$ in $\Gamma^+$.
The image of the workbench's norm-one order is
$\operatorname{Ad}(\mathcal O_{108}^{1})$, of index $6$ in this spinor
kernel and index $12$ in $\Gamma^+$.
\begin{proof}
The Hilbert symbol $(3,-1)_3$ is $(-1/3)=-1$, and the odd-unit formula
at $2$ gives $(3,-1)_2=(-1)^{((3-1)/2)((-1-1)/2)}=-1$.
At all other finite primes the symbol is $+1$, as at infinity.
Thus the algebra is an indefinite division algebra of discriminant $6$.
Direct multiplication checks that $\mathcal O_6$ is an order and that
its trace-pairing determinant is $-36$; it is maximal since its reduced
discriminant equals the algebra discriminant. The other discriminants
follow either from the trace pairings or from lattice indices $2$ and
$18$ in $\mathcal O_6$.

For comparison with the original presentation set
$I=3i+3j$, $J=i+3j$, $K=6k$. Then $I^2=18$, $J^2=-6$, $IJ=K=-JI$,
and the original basis $(1,(I+J)/2,(I-J)/2,K/2)$ gives exactly
$\mathcal O_{108}$. The map $V$ sends the parity lattice $a\equiv q\pmod2$
onto its trace-zero lattice, but the adjoint norm-one action preserves the
whole diagonal lattice as well. This last assertion, absent from the
workbench bridge, follows by direct multiplication. More generally, for
$g=r+si+cj+dk\in\mathcal O_{12}^1$, the only possible denominator in its
adjoint matrix in the $(a,q,n)$ coordinates is $3$, in a term $rc/3$.
The norm equation implies $r^2+c^2\equiv1\pmod3$, hence $3\mid rc$.
Thus $\operatorname{Ad}(\mathcal O_{12}^1)\subseteq\Gamma^+$.

At $2$, the quotient of the maximal order by its radical is $\mathbb F_4$,
and $\mathcal O_{12}$ is the preimage of $\mathbb F_2$. The norm-one
unit index is $3$: the unit $(1+i+3j+k)/2$ has norm $1$ and its reduction
generates $\mathbb F_4^\times$. At other primes the orders coincide.
Therefore $[\mathcal O_6^1:\mathcal O_{12}^1]=3$.
Voight's Theorem 39.1.2, in the opened version 0.9.14 of
\emph{Quaternion algebras}, states that a maximal order in an indefinite
rational quaternion algebra of discriminant $D$ has projective norm-one
covolume $\frac\pi3\prod_{p\mid D}(p-1)$ \cite{Voight2018}.
Here the hypotheses have just been verified and $D=6$, giving $2\pi/3$.
Thus the adjoint image of $\mathcal O_{12}^1$ has area $2\pi$ and index
$2$ in the already identified group of area $\pi$.

The quaternion $3+k$ has norm $6$ and adjoint matrix $T^{-1}$.
Its spinor class is the non-square $6$, whereas norm-one elements have
trivial class. Hence the index-two subgroup just computed is exactly
the spinor kernel.
Finally $\mathcal O_{108}^1$ is the subgroup of $\mathcal O_{12}^1$
with $c\equiv d\equiv0\pmod3$. The norm-one reductions modulo $3$
have $r^2+c^2=1$ and arbitrary $s,d$, hence size $36$; the indicated
subgroup has size $6$. Surjectivity follows from the reductions of the
actual norm-one units $j$, $2+i$, and $2+k$: the latter two generate the
two additive directions and $j$ generates the four residue units.
This proves index $6$, including the local conditions at both ramified
primes. The exact matrix calculations are reproduced in
\path{code/A_quaternion.py}.
\end{proof}

\subsection{Absence of exceptional spectrum}
\paragraph{Statement A.2.5 (Proved).}
Every compact hyperbolic orbifold of signature $(0;2,2,2,2,2)$,
independently of its moduli, has Cheeger constant at least $1$.
On $M=\Gamma^+\backslash\mathbb H_F$, every nonconstant Laplace
eigenfunction has eigenvalue strictly greater than $1/4$. In particular no
exceptional spectral term obstructs the exponent $2/3$ in either
the unrestricted or the arc-restricted problem.
\begin{proof}
We prove the stronger geometric fact that every compact hyperbolic
orbifold of signature $(0;2,2,2,2,2)$ has Cheeger constant at least $1$.
Here the Cheeger constant is the infimum of
$\operatorname{length}(\partial U)/\operatorname{area}(U)$ over
smooth domains of area at most half the total area. Boundaries can be
taken disjoint from the finite set of cone points, by approximation.

First let $D$ be a disk with a smooth boundary containing at most two
cone points. Its orbifold fundamental group maps into the ambient
Fuchsian group with image either trivial, cyclic of order two, or
generated by two half-turns about distinct points. In the last case
the image is infinite dihedral: the product of the half-turns is a
translation by twice the distance between their centres. The disk
lifts injectively and isometrically to the corresponding complete
orbifold cover. The injectivity of the lift follows immediately from
the identity obtained after composition with the covering map; no
claim about an embedded developed image in $\mathbb H_F$ is needed.

Each of these three complete covers has a smooth orbifold vector
field $V$ with $|V|\leq1$ and $\operatorname{div}V=1$.
For the trivial and finite elliptic cases use polar coordinates about
a fixed point and
$V=\tanh(r/2)\partial_r$ in the metric
$dr^2+\sinh^2r\,d\phi^2$; its divergence is
$1/(\cosh r+1)+\cosh r/(\cosh r+1)=1$.
For the infinite dihedral case use Fermi coordinates about the common
axis, with metric $dr^2+\cosh^2r\,ds^2$, and
$V=\tanh r\,\partial_r$. The field is invariant under translations
and under $(r,s)\mapsto(-r,-s)$, and its divergence is
$\operatorname{sech}^2r+\tanh^2r=1$.
The divergence theorem, or its application after deleting disks of
radius tending to zero about the cone points, now gives
\begin{equation}\label{A:disk-isoperimetry}
 \operatorname{area}(D)\leq\operatorname{length}(\partial D).
\end{equation}
For a nonempty disk the inequality is strict, since its lifted boundary
is compact and the displayed vector field has $|V|<1$ at every finite
point. The later argument only requires the weak inequality for the
Cheeger infimum, but the strict version will exclude eigenvalue $1/4$.


For completeness we pass from disks to arbitrary domains, a step
which matters for the Cheeger constant. The disjoint boundary circles
of $U$ cut the underlying sphere into components; their dual graph is
a tree. Give each vertex its number of cone points, with total weight
$5$. Orient each edge toward the side with at least three cone points.
There is a unique sink: two different sinks would require disjoint
subtrees each containing at least three of the five points. Every
subtree separated from the sink by one edge is a disk with at most
two cone points, so its area is bounded by the length of that edge's
boundary circle, by \eqref{A:disk-isoperimetry}.
If the sink is outside $U$, choose for each connected cluster of
vertices belonging to $U$ its boundary edge nearest the sink. The
associated disk contains the entire cluster. The chosen edges are
distinct, so their inequalities give
$\operatorname{length}(\partial U)\geq\operatorname{area}(U)$.
If the sink is inside $U$, apply the same argument to the complement:
$\operatorname{length}(\partial U)\geq
\operatorname{area}(M\setminus U)\geq\operatorname{area}(U)$.
Thus $h(M)\geq1$.

We also include the analytic implication to specify the convention
and to avoid a hidden orbifold hypothesis. For a smooth mean-zero
real function $f$, choose a median $m$ and put
$f_+=(f-m)_+$, $f_-=(m-f)_+$. Each support has area at most half
the total area. The coarea formula and the preceding isoperimetric
inequality yield
\[
 \int f_\pm^2\leq\int|\nabla(f_\pm^2)|
 \leq2\Big(\int f_\pm^2\int|\nabla f_\pm|^2\Big)^{1/2}.
\]
Adding the resulting inequalities gives
\[
 \int|\nabla f|^2\geq\frac14\int(f-m)^2
 \geq\frac14\int f^2.
\]
Orbifold charts justify coarea and integration by parts by lifting
to disks; isolated cone points have area zero. The Rayleigh
principle gives the weak spectral bound. For a nonconstant
eigenfunction the positive or negative part is nonzero, and its regular
positive level sets have positive area for a set of levels of positive
measure. The strict disk inequality, and the same tree argument, make
coarea strict on those levels. Thus its Rayleigh quotient is strictly
greater than $1/4$. We claim no uniform strict margin over all moduli.
Statement A.2.1
verified that our actual group has exactly this signature.
\end{proof}

\subsection{The unrestricted relative count}
\paragraph{Statement A.3.1 (Known theorem applied with complete bookkeeping).}
For $X\geq2$,
\begin{equation}\label{A:unrestricted-asymptotic}
 S(X)=\sum_{1\leq z\leq X}\rho(3z-1)\rho(3z+1)
 =\frac{9L_0^2}{2\pi^2}X+O(X^{2/3}).
\end{equation}
The implied constant is fixed for the displayed lattice and region.
\begin{proof}
We use Theorem 1.2 of Lekkas--Petridis and its explicitly proved
two-geodesic extension, Remark 6 \cite{LekkasPetridis2026}.
Their hypotheses are a cocompact torsion-free Fuchsian group and
primitive cyclic geodesic stabilizers. Their conclusion is
\[
 N(Y;l,l')=\sum_{1/2<s_j\leq1}\frac2\pi\gamma_1(s_j)
  \Big(\int_lu_j\,ds\Big)\Big(\int_{l'}u_j\,ds\Big)Y^{s_j}
  +O(Y^{2/3}),
\]
where the $u_j$ are orthonormal Laplace eigenfunctions,
$\lambda_j=s_j(1-s_j)$, and
\[
 \gamma_1(s)=\frac\pi2
 \frac{\Gamma(s+1/2)\Gamma((s+1)/2)}{\Gamma(s/2)^3}
 \frac4{s(2s-1)},\qquad \gamma_1(1)=1.
\]
Intersecting axes contribute only finitely many double cosets.
The extension in Remark 6 replaces squared periods by products;
its proof uses the same transform and Cauchy--Schwarz for the two
period sums. We now verify the torsion removal and its exact effect,
instead of applying the torsion-free statement directly to $M$.

Set $K=\ker(\Gamma^+\to\operatorname{GL}_3(\mathbb F_5))$ and
$d=[\Gamma^+:K]$. It is normal of finite index and cocompact.
It is torsion-free: a nonidentity matrix congruent to $1$ modulo $5$
can be written $1+5^eB$ with $e\geq1$ and $B\not\equiv0\pmod5$.
If it has finite order, a power has prime order $p$. For $p\ne5$
the binomial expansion modulo $5^{e+1}$ contradicts order $p$;
for $p=5$ expansion modulo $5^{e+2}$ gives the same contradiction.
Put $H_i'=H_i\cap K$, $m_i=[H_i:H_i']$. These are the full
oriented cyclic stabilizers in $K$ and their primitive lengths are
$m_i l_i$. The values from Statement A.2.2 are
$l_0=L_0$, $l_1=2L_0$, $l_2=L_0$.

For $i=1,2$ no translate of the target axis equals the source axis:
such equality would make an integral norm-$-2$ normal proportional
to $(0,0,1)$, hence equal to $(0,0,\pm1/3)$, which is impossible.
Consequently $H_0\cap\gamma H_i\gamma^{-1}$ is trivial for every
$\gamma$. The natural map
\[
 H_0'\backslash\Gamma^+/H_i'\longrightarrow
 H_0\backslash\Gamma^+/H_i
\]
therefore has exactly $m_0m_i$ preimages of every double coset:
they are represented by the $m_0$ left and $m_i$ right residue
classes, and an equality between two of them would give a
nonidentity element of that intersection. Distance is unchanged.

Decompose $\Gamma^+$ into its $d$ cosets $K\gamma_r$.
The part in $K\gamma_r$ is the pair count in $K$ for $H_0'$ and
$\gamma_rH_i'\gamma_r^{-1}$. All these stabilizers are primitive.
There are no asymptotic pairs of distinct axes: two hyperbolic
elements of a discrete cocompact Fuchsian group sharing one boundary
fixed point share both, and then their axes coincide. To see the
claim directly, conjugate the common endpoint to infinity; distinct
second endpoints and conjugation by powers of one dilation produce
nonidentity group elements tending to the identity. Thus the
regularity condition $abcd\ne0$ in the cited construction holds
outside the finite intersection terms.

Apply the displayed theorem to each of the finitely many pairs in
$K\backslash\mathbb H_F$. Extra small eigenvalues of this cover do
not survive the sum. Indeed the deck group $\Gamma^+/K$ acts
unitarily on every Laplace eigenspace. The sum of the target periods
over $\gamma_r$ is $d$ times the period of the orthogonal projection
onto the deck-invariant subspace. We may choose an orthonormal
eigenbasis respecting that subspace. All other terms cancel exactly.
The invariant eigenfunctions are $d^{-1/2}$ times normalized
eigenfunctions on $M$. Their source and target periods acquire
factors $m_0/\sqrt d$ and $m_i/\sqrt d$, respectively. Summing $d$
cosets and dividing by $m_0m_i$ cancels these factors. By Statement
A.2.5 only the constant eigenfunction has $1/2<s_j\leq1$.
We obtain, for each $i$,
\[
 N_i(Y)=\frac{2l_0l_i}{\pi\operatorname{area}(M)}Y
       +O(Y^{2/3}).
\]
Since $\operatorname{area}(M)=\pi$, summing $i=1,2$, substituting
$Y=3X$, and using the exact identity
$N_1(3X)+N_2(3X)=4S(X)+2$ proves
\eqref{A:unrestricted-asymptotic}.

The theorem is written with a strict cutoff, whereas our count uses
a weak cutoff. Their difference at a height $z$ is at most
$4\rho(m_z)$. The elementary estimate $d(n)\ll_\epsilon n^\epsilon$
therefore makes this $O_\epsilon(X^\epsilon)$, which is absorbed.
For clarity that elementary estimate follows by separating the
finitely many primes $p<2^{1/\epsilon}$; for the remaining primes
$e+1\leq2^e\leq p^{\epsilon e}$, and the finitely many maxima of
$(e+1)p^{-\epsilon e}$ are finite.
\end{proof}

\paragraph{Statement A.3.2 (Proved reformulation and literature comparison).}
Writing $n=3z-1$ gives the exact shifted norm-coefficient sum
\[
 S(X)=\sum_{\substack{1\leq n\leq3X-1\\n\equiv2\pmod3}}
       (1*\chi_{12})(n)(1*\chi_{12})(n+2).
\]
These are the coefficients counting ideals in $\mathbb Q(\sqrt3)$:
$\zeta_{\mathbb Q(\sqrt3)}(s)=\zeta(s)L(s,\chi_{12})$.
This specifies the binary additive-divisor interpretation precisely:
the additive shift $2$ is in the integer norm variable, with the
progression $n\equiv2\pmod3$. It is not a sum over elements of the
quadratic field modulo an unstated unit convention.
For comparison, Motohashi's Theorem 1 and Corollary 1
\cite{Motohashi1994} give, for every fixed positive shift $h$,
\[
 \sum_{n\leq Y}d(n)d(n+h)=Y P_{2,h}(\log Y)
       +O_{h,\epsilon}(Y^{2/3+\epsilon}),
\]
where the quadratic polynomial has leading coefficient
$6\sigma_{-1}(h)/\pi^2$. In our case the main term is linear without a
logarithmic polynomial. The distinction is consistent with the simple
pole of the Dedekind zeta function, as opposed to the double pole of
$\zeta(s)^2$. The cited classical theorem concerns $d(n)$, so it is a
comparison and is not used as a theorem about the character-twisted
coefficients or the sharp ideal arc. Statement A.3.1 proves our formula
from the relative hyperbolic count instead.

\subsection{The sharp arc: an explicit power saving and the uniform transform}
\paragraph{Statement A.4.1 (Proved; weaker than the requested exponent).}
For every $\epsilon>0$,
\begin{equation}\label{A:sharp-7-8}
 R(X)=\frac{9L_0b_0}{2\pi^2}X+O_\epsilon(X^{7/8+\epsilon}).
\end{equation}
In particular this proves the proposed leading constant for the actual
sharp arc. It does not prove the requested exponent $2/3$.
\begin{proof}
Here is a direct application of the general lattice-counting theorem of
Gorodnik--Nevo \cite{GorodnikNevo2012}, with its domain regularity verified
for this arc. Work first in $G=\operatorname{PSL}_2(\mathbb R)$, identified
with the connected isometry group of $\mathbb H_F$. Conjugate the source
axis to the imaginary axis and put
\[
 A_u=\begin{pmatrix}e^{u/2}&0\\0&e^{-u/2}\end{pmatrix},\qquad
 B_r=\begin{pmatrix}\cosh(r/2)&\sinh(r/2)\\
                  \sinh(r/2)&\cosh(r/2)\end{pmatrix}.
\]
For $g=A_uB_rA_v=\begin{pmatrix}a&b\\c&d\end{pmatrix}$,
\[
 ad+bc=\cosh r,\qquad
 u=\frac12\log\left|\frac{ab}{cd}\right|,\qquad
 v=\frac12\log\left|\frac{ac}{bd}\right|.
\]
The parameter $u$ is arclength of the perpendicular foot on the source
axis, with the choice of origin and orientation in Statement A.2.3.
There are four regular cells, according to the sign of $ad+bc$ and the
side/orientation of the perpendicular; each is obtained from the displayed
cell by fixed matrices reversing axes and/or replacing $r$ by $-r$.
Their absolute Haar densities agree. Direct differentiation of
$g^{-1}dg$ gives
\begin{equation}\label{A:Haar-ABA}
 dg=\kappa\sinh r\,du\,dr\,dv\qquad(r>0),
\end{equation}
for a positive normalization constant $\kappa$. The exact symbolic check
is in \path{code/A_relative_coordinates.py}.

Choose $\tau_i\in G$ carrying the standard oriented normal to the
normal of the $i$th target geodesic, for $i=1,2$. Then
$H_i=\tau_i\langle A_{l_i}\rangle\tau_i^{-1}$. An element of
$H_0\backslash\Gamma^+/H_i$ has a unique representative in the coset
$\Gamma^+\tau_i$ with
\[
 0\leq u<L_0,\qquad 0\leq v<l_i,
\]
once its regular cell is specified. Indeed left and right translations
change $u$ and $v$ by the indicated primitive lengths, and their
intersection is trivial, as proved in Statement A.3.1. A single one of
the four cells is the positive-$s$, positive-$z$ ideal count. In that
cell the sharp arc is exactly $0\leq u\leq b_0$.
The lower bound $r\geq1$ removes no positive integral height, since
$\cosh r=3z\geq3$ and $\operatorname{arcosh}3>1$.
Thus, apart from the explicitly treated endpoints, its count is the
number of coset points in
\[
 D_i(X;I)=\{A_u B_r A_v:u\in I,\ 0\leq v<l_i,
              \ 1\leq r\leq\operatorname{arcosh}(3X)\}
\]
in the appropriate regular cell, where $I=[0,b_0]$ or $[0,L_0)$.
Changing which endpoint of a fundamental interval is included selects
one representative and introduces no multiplicity.

We verify Lipschitz well-roundedness, rather than assuming a sector
regularity theorem. On each such domain all four entries have absolute
value between fixed positive multiples of $e^{r/2}$: the ranges of
$u,v$ are fixed compact intervals and $r\geq1$. Left or right
multiplication by a matrix $1+O(\delta)$ changes each entry relatively
by $O(\delta)$, uniformly in $r$. The signs therefore remain fixed, the
two logarithmic formulas change $u,v$ by $O(\delta)$, and the formula
for $\cosh r$ changes $r$ by $O(\delta)$, since $\coth r$ is bounded
for $r\geq1$. The same estimates applied to inverse perturbations give
an inner comparison. Hence the outer and inner $\delta$-perturbations
are bounded by coordinate rectangles whose six faces have moved by
$O(\delta)$. Formula \eqref{A:Haar-ABA} shows
\[
 \operatorname{vol}(D_i^+(\delta))-
 \operatorname{vol}(D_i^-(\delta))
       =O(\delta\operatorname{vol}(D_i)),\qquad
 \operatorname{vol}(D_i)\asymp X.
\]
This is precisely Definition 1.1 of the cited paper with exponent
$a=1$. Small identity balls have upper local dimension $\rho=3$.

The representation on $L^2_0(\Gamma^+\backslash G)$ is $L^{2+}$.
For completeness, the unitary dual classification
\cite[Theorem (2.3)]{Moriyama1999} consists of the trivial representation,
principal series, discrete series and their limits, and complementary
series. Every complementary series has a $K$-fixed vector and corresponds
to a scalar Laplace eigenvalue in $(0,1/4)$, excluded by Statement A.2.5.
The other nontrivial representations are tempered: their $K$-finite
matrix coefficients are in $L^{2+}$ (in polar coordinates the principal
series bound is a constant times $(1+r)e^{-r/2}$ and Haar density is
$O(e^r)$; discrete series are square integrable, with their limits
satisfying the same $L^{2+}$ bound). Finite sums of such vectors are dense.
The constant representation is removed, and connectedness of the quotient
ensures there are no other invariant vectors.

To satisfy the algebraic simply-connected hypothesis in Gorodnik--Nevo,
lift the group, lattice and domains to $\operatorname{SL}_2(\mathbb R)$.
This is a real isotropic simply connected semisimple algebraic group;
the twofold central lift changes point counts and volumes by the same
factor. It does not change the representation on the quotient.
Theorem 4.5 and equation (3.1) of the cited paper, with $p=2$ and
$n_e(p)=1$, give the stable operator bound
$E(X)\ll_\eta\operatorname{vol}(D_i)^{-1/2+\eta}$.
The lattice is discrete and cocompact, so the identity and the finitely
many translates $\tau_i$ have positive injectivity radii. Theorem 1.9
and Corollary 1.11, with $a=1,\rho=3$, now give
\[
 \#(\Gamma^+\tau_i\cap D_i(X;I))
 =\frac{\operatorname{vol}(D_i(X;I))}
        {\operatorname{vol}(\Gamma^+\backslash G)}
    +O_\epsilon(X^{7/8+\epsilon}).
\]
This checks both the translated-coset and sharp-boundary hypotheses.

It remains to normalize the main coefficient. By
\eqref{A:Haar-ABA}, restriction from $[0,L_0)$ to $[0,b_0]$ multiplies
volume by $b_0/L_0$. The four regular cells have equal volume. Their
full count is $4S(X)+O(1)$ by Statement A.2.3, whose leading coefficient
was proved in Statement A.3.1. Consequently the positive-$s$, positive-$z$
cell has leading coefficient $9L_0^2/(2\pi^2)$, and the arc has coefficient
$9L_0b_0/(2\pi^2)$, as asserted.
The lower endpoint consists of the Pell family $t=0$, hence contributes
$O(\log X)$ if its convention is changed. The upper endpoint is empty:
$288s^2=289(18z^2-2)$ would imply $17\mid s$ and then
$144(s/17)^2=9z^2-1$, impossible modulo $9$.
Finally Statement A.1.3 changes the arc count by $O(\log X)$.
\end{proof}


\paragraph{Statement A.4.2 (Proved reduction; its needed uniform estimate remains undecided).}
For each integer $\ell$ define a unitary character on the nonzero ideals
of $\mathbb Z[\sqrt3]$ by
\[
 \psi_\ell((\alpha))=
 \left|\frac{\alpha}{\alpha'}\right|^{i\pi\ell/L_0},
 \qquad a_\ell(m)=\sum_{N\mathfrak a=m}\psi_\ell(\mathfrak a),
\]
where the prime denotes the nontrivial real embedding. Then
\begin{equation}\label{A:twisted-shift}
 T_\ell(X):=\sum_{z\leq X}\sum_{N\mathfrak a=m_z}
 e^{2\pi i\ell\theta(\mathfrak a)}
 =(-1)^\ell\sum_{z\leq X}a_\ell(3z-1)a_\ell(3z+1).
\end{equation}
For every integer $H\geq1$ the actual sharp-arc error satisfies
\begin{equation}\label{A:ET-reduction}
 |R(X)-C_0X|\ll X^{2/3}+\frac XH+
       \sum_{1\leq\ell\leq H}\frac{|T_\ell(X)|}{\ell}
       +\log(2X).
\end{equation}
In particular, the uniform bound
\begin{equation}\label{A:needed-twist}
 |T_\ell(X)|\ll_\epsilon X^{2/3+\epsilon}
           +X^{1/2+\epsilon}\ell^{1/2}
 \quad(1\leq\ell\leq X^{1/3})
\end{equation}
would prove the requested exponent $2/3+\epsilon$.
We have not proved \eqref{A:needed-twist}.
\begin{proof}
Changing the generator by $\pm\varepsilon_0^k$ changes the logarithm
of the absolute embedding ratio by $2kL_0$, and so leaves
$\psi_\ell$ unchanged. The character is multiplicative. On the
positive-norm generator of $m_z$ it is exactly
$e^{2\pi i\ell\theta}$. For the unique prime ideal above $2$,
$|(1+\sqrt3)/(1-\sqrt3)|=\varepsilon_0$, so its character is
$(-1)^\ell$. Away from $2$, the integers $3z-1$ and $3z+1$ are
coprime. At $2$, where they can share powers, there is just one prime
ideal, and its local coefficient is geometric. Combining these facts
proves \eqref{A:twisted-shift}, including the ramified factor.
Conjugation pairs inverse character values, so $a_\ell(m)$ is real.

Apply the Erd\H{o}s--Tur\'an inequality in the form of
\cite[equation (8.2)]{CarneiroVaaler2010} to the finite list of the
$S(X)$ ideal angles. That equation bounds the translate of the sawtooth
sum; the characteristic function of an interval minus its length is
the difference of two translates of that sawtooth function, away from
its endpoints. Thus the arc discrepancy is at most a fixed constant
times $S(X)/H+\sum_{\ell\leq H}|T_\ell(X)|/\ell$, plus endpoint
multiplicities. The upper endpoint is empty by the calculation in
Statement A.4.1. At the lower endpoint, $t=0$ gives
$s^2-2(3z)^2=-2$. Reduction in $\mathbb Z[\sqrt2]$ by its norm-one
unit $3+2\sqrt2$ puts a positive generator in
$[\sqrt2,(3+2\sqrt2)\sqrt2)$, where the equation of norm $-2$
has the sole seed $(s,3z)=(0,1)$. Thus the total number with
$z\leq X$ is $O(\log(2X))$; imposing divisibility of the second
coordinate by $3$ can only decrease it.
Statement A.3.1 gives $S(X)=cX+O(X^{2/3})$ and $S(X)\ll X$.
Statement A.1.3 now proves \eqref{A:ET-reduction}.
Under \eqref{A:needed-twist}, take $H=\lfloor X^{1/3}\rfloor$;
the two sums cost $O_\epsilon(X^{2/3+\epsilon}\log X)$ and
$O_\epsilon(X^{1/2+\epsilon}H^{1/2})$, respectively. Relabeling
$\epsilon$ proves the asserted conditional implication.
\end{proof}

\paragraph{Uniformity audit and current obstruction.}
A bound for each fixed $\ell$ with an unspecified constant depending on
$\ell$ cannot be inserted in \eqref{A:ET-reduction} with $H$ growing
with $X$. More quantitatively, a bound
$T_\ell(X)\ll_\epsilon \ell^A X^{\beta+\epsilon}$, $A>0$, leads by this
argument to exponent $(A+\beta)/(A+1)$, after balancing $X/H$ against
$X^\beta H^A$. Thus even a fixed-mode exponent $2/3$ with polynomial
loss in the mode does not establish the sharp-arc exponent $2/3$.
Theorems 1.2 and 1.5 of Lekkas--Petridis integrate the entire source
geodesic. Their Remark 6 changes the target geodesic, but does not
supply these growing-mode estimates. Voskou's four orientation cases
\cite{Voskou2024} are discrete sign cases, not intervals in the foot
parameter. The fixed compact angular modes in
\cite[Theorem 1.1]{BourgainKontorovichSarnak2010} likewise do not
supply the required uniform hyperbolic-foot estimate.
The direct domain method of Statement A.4.1 proves $7/8+\epsilon$.
Statement A.4.4 supplies a uniform averaged substitute by Mellin
inversion and proves the requested $2/3+\epsilon$ exponent. The
stronger individual-mode bound \eqref{A:needed-twist} is not needed
for that proof and is not asserted here.


\paragraph{Statement A.4.3 (Proved uniform spectral coefficient estimate).}
Let $M$ be a fixed compact hyperbolic surface, let $l_0,l_1$ be fixed
closed geodesics of lengths $L_0,L_1$, and let $u_j$ be an orthonormal
basis of scalar Laplace eigenfunctions with
$\lambda_j=1/4+t_j^2$, $t_j\geq0$. Set $s_j=1/2+it_j$ and
\[
 P_{j,\ell}=L_0^{-1/2}\int_{l_0}u_j(u)e^{2\pi i\ell u/L_0}\,du,
 \quad P_j'=L_1^{-1/2}\int_{l_1}u_j\,du,\quad
 \kappa_\ell=2\pi\ell/L_0,
\]
\[
 D_s(\kappa)=
 \frac{\Gamma((s+i\kappa)/2)\Gamma((s-i\kappa)/2)}{\Gamma(s)},
 \qquad
 Q_{j,H}=\sum_{|\ell|\leq H}b_\ell
             \frac{D_{s_j}(0)}{D_{s_j}(\kappa_\ell)}P_{j,\ell},
\]
where $|b_\ell|\leq B/(1+|\ell|)$ and $H\geq1$. Then, for $T\geq1$,
\begin{equation}\label{A:uniform-arc-period}
 \sum_{T\leq t_j<2T}|Q_{j,H}|^2
       \ll B^2T\log^2(2H),\qquad
 \sum_{T\leq t_j<2T}|P_j'Q_{j,H}|
       \ll BT\log(2H).
\end{equation}
The constants depend on the surface and geodesics, but not on $T,H$
or the displayed coefficients. Thus the angular sum itself can be
controlled without a power of $H$ in this normalized spectral expression.
The uniform inversion that turns this coefficient bound into a sharp-arc
count is proved in Statement A.4.4 below.
\begin{proof}
We first record an elementary broad-band restriction estimate, uniformly
in the real frequency $\kappa$:
\begin{equation}\label{A:broad-band-period}
 \sum_{t_j\leq2T}\left|\int_{l_0}u_j(u)e^{i\kappa u}\,du\right|^2
       \ll T.
\end{equation}
Indeed, in finitely many Fermi charts along the parametrized geodesic,
the one-dimensional fundamental theorem of calculus and Cauchy--Schwarz
give, for $0<\delta\ll1$,
\[
 \|f|_{l_0}\|_2^2\ll
 \delta^{-1}\|f\|_{L^2(M)}^2+
 \delta\|\nabla f\|_{L^2(M)}^2.
\]
The finite number of charts accounts for any self-intersections.
For a linear combination of eigenfunctions with $t_j\leq2T$,
$\|\nabla f\|_2^2\leq(4T^2+1/4)\|f\|_2^2$.
Take $\delta=T^{-1}$ (adjust a fixed constant for small $T$), and
apply the adjoint restriction operator to $e^{i\kappa u}$, whose
$L^2(l_0)$ norm is fixed. This proves \eqref{A:broad-band-period}.
The same estimate applies to $P_j'$.
We also use the elementary bound $\#\{t_j\leq2T\}\ll T^2$.
One proof partitions a fixed finite collection of coordinate patches
into $O(T^2)$ cells of diameter $c/T$ and uniformly bounded shape.
If the spectral space had larger dimension, a nonzero member would
have integral zero on every cell. The cell Poincar\'e inequalities would
give $\|f\|_2^2\leq C(c/T)^2\|\nabla f\|_2^2$, contradicting the
spectral bound when the fixed $c$ is sufficiently small.

Reznikov's mean-value estimate
\cite[equation (1.13), proved in Sections 2.3 and 3.1--3.2]{Reznikov2015}
states that the normalized hyperbolic coefficients $a_{j,\ell}$ satisfy
\begin{equation}\label{A:Reznikov-mean}
 \sum_{|\ell|\leq N}|a_{j,\ell}|^2\ll\max(N,1+t_j),
 \qquad P_{j,\ell}=a_{j,\ell}D_{s_j}(\kappa_\ell).
\end{equation}
A fixed scalar from the model functional in his equation (2.3) has
been absorbed into $a_{j,\ell}$. Explicitly, put his representation
parameter equal to $-2it_j$ and his diagonal-character parameter equal
to $2i\kappa_\ell$. His gamma arguments
$(1-\lambda_{\mathrm{Rep}}\pm s_{\mathrm{char}})/4$ then become
$(s_j\pm i\kappa_\ell)/2$, and his denominator argument
$(1-\lambda_{\mathrm{Rep}})/2$ is $s_j$. The scalar $2/\pi$
is independent of both parameters. The hypotheses are precisely a fixed
closed geodesic on a compact hyperbolic surface and an $L^2$-normalized
principal-series eigenfunction. His argument includes both components
of the diagonal group; retaining the component that gives scalar
periods only decreases its nonnegative sum. The assertion is uniform
in the eigenfunction as well as in $N$.

Split the angular modes at $|\kappa_\ell|=T/2$.
For $t_j\in[T,2T]$ and the lower modes, Stirling's formula in the
fixed vertical strip $\operatorname{Re}(s_j/2)=1/4$ gives
\[
 \left|\frac{D_{s_j}(0)}{D_{s_j}(\kappa_\ell)}\right|\ll1.
\]
For bounded $T$ this follows from continuity and the absence of gamma
zeros or poles; for large $T$ both $t_j\pm\kappa_\ell$ are comparable
to $T$ and the exponential factors cancel exactly.
Minkowski's inequality in the $j$ variable and
\eqref{A:broad-band-period} bound the square root of the lower-mode
sum by
\[
 \sum_{|\kappa_\ell|\leq T/2}|b_\ell|
 \left(\sum_{T\leq t_j<2T}|P_{j,\ell}|^2\right)^{1/2}
       \ll B\sqrt T\log(2H).
\]

For the higher modes divide the integers into dyadic blocks
$N\leq|\ell|<2N$, with $N\gg T$ up to a fixed geodesic-length factor.
On such a block \eqref{A:Reznikov-mean} and Cauchy--Schwarz give
\[
 \left|\sum_{N\leq|\ell|<2N}b_\ell a_{j,\ell}\right|
 \leq\left(\sum|a_{j,\ell}|^2\right)^{1/2}
       \left(\sum|b_\ell|^2\right)^{1/2}\ll B.
\]
Stirling's formula also gives $|D_{s_j}(0)|\ll T^{-1/2}$.
Since the gamma factors cancel exactly in the definition of $Q$,
the higher-mode contribution to each $Q_{j,H}$ is at most
$BT^{-1/2}\log(2H)$. There are $O(T^2)$ eigenvalues in the band,
so its squared sum is $O(B^2T\log^2(2H))$.
Combining the two ranges proves the first estimate in
\eqref{A:uniform-arc-period}; Cauchy--Schwarz with
$\sum|P_j'|^2\ll T$ proves the second.
\end{proof}

The estimate applies after torsion removal to the fixed cover used in
Statement A.3.1. The original angular frequencies are a fixed-index
subsequence of the cover's frequencies, and the constants in
\eqref{A:Reznikov-mean} change only by that fixed index. The next statement derives the corresponding uniform radial transform
and proves its inversion estimates explicitly.

\paragraph{Statement A.4.4 (Proved sharp-arc exponent).}
For every $\epsilon>0$, the actual count in the question satisfies
\begin{equation}\label{A:sharp-two-thirds}
 R(X)=C_0X+O_\epsilon(X^{2/3+\epsilon}).
\end{equation}
The proof uses an angularly uniform Mellin transform. In particular, it
does not deduce an arc estimate from a count over an entire geodesic.

\begin{proof}
We give the transform and all uniform estimates used to invert it.
Put $A_z=(9z^2-1)^{1/2}$ and $\kappa_\ell=2\pi\ell/L_0$.
In this proof periods are unnormalized; this changes the estimates of
Statement A.4.3 only by fixed geodesic-length factors. Choose real
orthonormal eigenfunctions $u_j$ on $\Gamma^+\backslash\mathbb H$ and put
\[
 p_{j,\ell}=\int_{l_0}u_j(u)e^{i\kappa_\ell u}\,du,
 \qquad p_j'=\sum_{i=1}^2\int_{l_i}u_j\,du.
\]
Apart from the constant eigenfunction, write $s_j=1/2+it_j$, with
$t_j>0$, using Statement A.2.5. The constants in the following estimates
may depend on this fixed orbifold and its three geodesics. The spectral
estimates on its torsion-free cover apply to these eigenfunctions:
their lifts are normalized by the square root of the cover degree,
and a period on a lifted geodesic is the original period multiplied by
its fixed covering degree. The angular frequencies form the corresponding
fixed-index subsequence. Thus no spectral gap is assumed on the cover.

\emph{1. An exact Mellin identity.}
For $\operatorname{Re}s>1$ form the absolutely convergent series
\[
 E_s(p)=\sum_{i=1}^2\sum_{H_i\backslash\Gamma^+}
       \cosh\bigl(d(p,\gamma^{-1}l_i)\bigr)^{-s}.
\]
Absolute local uniform convergence follows by placing fixed disjoint
balls around coset representatives in the hyperbolic cylinder:
the number at distance at most $r$ is $O(e^r)$, since a cylinder
of fixed neck length has area $O(e^r)$ in that range. Its derivatives
also converge locally after increasing $\operatorname{Re}s$; the
identities initially obtained there continue to $\operatorname{Re}s>1$.
Set
\[
 K_s(\kappa)=\int_{\mathbb R}(\cosh u)^{-s}e^{i\kappa u}\,du
 =\frac{2^{s-1}\Gamma((s+i\kappa)/2)\Gamma((s-i\kappa)/2)}{\Gamma(s)},
\]
\[
 M_\ell(s)=K_s(\kappa_\ell)^{-1}
                 \int_{l_0}E_s(u)e^{i\kappa_\ell u}\,du.
\]
The beta integral proves the displayed formula for $K_s$.
For a disjoint target axis, the signed distance of a point on $l_0$
from that axis has absolute hyperbolic sine
$A\cosh(u-u_0)$, where $A=\sinh d(l_0,l)$ and $u_0$ is the
perpendicular foot. Consequently its contribution to $M_\ell$ is
\[
 e^{i\kappa_\ell u_0}A^{-s}F_{s,\kappa_\ell}(A),\quad
 F_{s,\kappa}(A)=
 {}_2F_1\left(\frac{s+i\kappa}{2},\frac{s-i\kappa}{2};
                    \frac{s+1}{2};-A^{-2}\right).
\]
To verify this identity without an unquoted transform formula, expand
$(1+A^{-2}\operatorname{sech}^2u)^{-s/2}$ and integrate termwise.
The coefficient is
\[
 \frac{(-1)^k(s/2)_k}{k!A^{2k}}
      \frac{K_{s+2k}(\kappa)}{K_s(\kappa)}
 =\frac{(-1)^k((s+i\kappa)/2)_k((s-i\kappa)/2)_k}
          {k!((s+1)/2)_k A^{2k}},
\]
by the gamma duplication formula. Here $A=A_z\geq\sqrt8$, so the
series and the original integrals converge absolutely.

There are four copies of each positive-height ideal angle: changing
the normal's sign and changing the sign of its $z$ coordinate leave
the foot coordinate unchanged modulo $L_0$. These are the four
regular cells already counted in Statement A.2.3. The two exceptional
double cosets with $z=0$ meet $l_0$ orthogonally at $u=L_0/2$.
Indeed $r_2=(1,1,0)$ is orthogonal to $e_z$ and its intersection
point has $\tanh u=1/\sqrt3$. Their individual contributions are
$e^{i\kappa_\ell L_0/2}=(-1)^\ell$. Thus exactly
\begin{equation}\label{A:Mellin-geometry}
 M_\ell(s)=2(-1)^\ell+4D_\ell(s)+4H_\ell(s),
\end{equation}
where
\[
 D_\ell(s)=\sum_{z\geq1}\sum_{N\mathfrak a=m_z}
          e^{2\pi i\ell\theta(\mathfrak a)}A_z^{-s},
\]
\[
 H_\ell(s)=\sum_{z\geq1}\sum_{N\mathfrak a=m_z}
 e^{2\pi i\ell\theta(\mathfrak a)}A_z^{-s}
                    (F_{s,\kappa_\ell}(A_z)-1).
\]
All sums in \eqref{A:Mellin-geometry} are absolutely convergent
in their initial half-plane. The half-turns reversing an axis do not
belong to its cyclic subgroup $H_i$; they therefore give the two
oriented normal cosets, as required for this factor of four.

Unfolding the scalar spectral coefficient of $E_s$ along $l_i$ gives
\[
 \begin{aligned}
 &\left(\int_{l_i}u_j\right)
 \int_{\mathbb R}(1+x^2)^{-s/2}
 {}_2F_1\left(\frac{s_j}{2},\frac{1-s_j}{2};\frac12;-x^2\right)dx
 \\
 &\qquad=\left(\int_{l_i}u_j\right)
 \frac{\sqrt\pi\Gamma((s+s_j-1)/2)\Gamma((s-s_j)/2)}{\Gamma(s/2)^2}.
 \end{aligned}
\]
Here $x=\sinh r$ in Fermi coordinates. Integrating the eigenfunction
along the full target geodesic solves
$\xi''+\tanh r\,\xi'+\lambda_j\xi=0$; its odd initial-data solution
integrates to zero over $r\in\mathbb R$, and its even solution is
the displayed hypergeometric function. For the beta evaluation put
$y=x^2/(1+x^2)$, apply the power-series identity
\[
 {}_2F_1(s_j/2,(1-s_j)/2;1/2;-x^2)
 =(1+x^2)^{-s_j/2}{}_2F_1(s_j/2,s_j/2;1/2;y),
\]
and integrate the series against
$y^{-1/2}(1-y)^{(s+s_j-3)/2}$. The resulting Gauss beta sum
is the stated gamma quotient. All beta integrals are absolutely
convergent for $\operatorname{Re}s>1$.
It follows that, for $\ell\ne0$,
\begin{equation}\label{A:Mellin-spectral}
 M_\ell(s)=\sum_{j\geq1}p_j'p_{j,\ell}G(s,t_j,\kappa_\ell),
\end{equation}
\[
 G(s,t,\kappa)=
 \frac{\Gamma((s+1)/2)\Gamma((s+s_j-1)/2)\Gamma((s-s_j)/2)}
 {\Gamma(s/2)\Gamma((s+i\kappa)/2)\Gamma((s-i\kappa)/2)}.
\]
The constant eigenfunction vanishes because its nonzero Fourier
period on $l_0$ is zero. For each fixed $\ell$ the series in
\eqref{A:Mellin-spectral} converges normally on compact subsets away
from its poles, since its terms decay exponentially as
$t_j\to\infty$. It therefore supplies the needed meromorphic
continuation. In $\operatorname{Re}s>0$ its only possible poles are
$s=s_j$ and $s=1-s_j$.

\emph{2. A uniform bound for the geometric correction.}
Fix $0<\sigma<1/2$, write $s=\sigma+i\tau$, and put
$U=1+|\tau|$. With $\delta=\operatorname{arsinh}(1/A)$ one has
\begin{equation}\label{A:F-spherical-integral}
 F_{s,\kappa}(A)=
 \frac{\Gamma((s+1)/2)}{\sqrt\pi\Gamma(s/2)}
 \int_{-1}^1
 \bigl(\cosh(2\delta)+x\sinh(2\delta)\bigr)^{-(s+i\kappa)/2}
 (1-x^2)^{s/2-1}\,dx.
\end{equation}
For completeness, this identity follows by differentiating under
the integral and integrating once by parts in $x$: both sides,
as analytic even functions of $\delta$ near zero, solve
\[
 f''+2s\coth(2\delta)f'+(s^2+\kappa^2)f=0,
 \qquad f(0)=1,\quad f'(0)=0.
\]
Uniqueness follows coefficient by coefficient, since
$\operatorname{Re}s>0$; continuation along positive real $\delta$
then proves \eqref{A:F-spherical-integral}. Equivalently, it is
the quadratic hypergeometric transformation with parameters
$a+b=s$, $c=(s+1)/2$.
The base in the integral is positive. Its complex power has modulus
equal to the power $-\sigma/2$, independently of both $\tau$ and
$\kappa$; the weight has modulus $(1-x^2)^{\sigma/2-1}$.
Stirling's formula for the prefactor gives
\[
 |F_{s,\kappa}(A)|\ll_\sigma U^{1/2}.
\]
At $\delta=0$ the first derivative of the integrand is an odd
multiple of $x$, whose integral is zero. Its second derivative is
bounded by a constant times $(1+|s|+|\kappa|)^2$ times the same
integrable absolute weight, uniformly for
$0\leq\delta\leq\operatorname{arsinh}(1/\sqrt8)$.
Taylor's formula with integral remainder therefore gives
\begin{equation}\label{A:F-uniform}
 |F_{s,\kappa}(A)-1|
 \ll_\sigma U^{1/2}
       \min\{1,(U+|\kappa|)^2/A^2\}.
\end{equation}
Statement A.3.1 implies
$\#\{(z,\mathfrak a):A_z\leq V\}\ll V$.
Partial summation in \eqref{A:F-uniform} now proves normal
absolute convergence of $H_\ell(s)$ for $\operatorname{Re}s>0$
and, on our chosen line,
\begin{equation}\label{A:H-uniform}
 |H_\ell(\sigma+i\tau)|
       \ll_\sigma U^{1/2}(U+|\kappa_\ell|)^{1-\sigma}.
\end{equation}
The same convergence argument continues it to any closed strip
in $\operatorname{Re}s>-1$ avoiding the possible hypergeometric
parameter poles, but only the positive strip is needed here.

\emph{3. The remaining Mellin contour, uniformly in the angular cutoff.}
Let $T,H\geq2$. Let $w=w_T$ be a smooth function, equal to one
near zero, supported in $[0,2]$, with its transition confined to
$[1-1/T,1+1/T]$, $0\leq w\leq1$, and
$|w^{(k)}|\ll_kT^k$. Its Mellin transform satisfies, on every
fixed positive vertical strip and for every $N$,
\begin{equation}\label{A:w-Mellin-bound}
 |\widehat w(\sigma+i\tau)|
 \ll_{\sigma,N}U^{-1}(1+U/T)^{-N}.
\end{equation}
This follows by one integration by parts, whose derivative has
bounded $L^1$ norm, and repeated further integrations by parts in
the logarithmic variable. Let $b_\ell$, $1\leq\ell\leq H$, be
arbitrary complex numbers with $|b_\ell|\leq1/\ell$.
We prove
\begin{equation}\label{A:contour-uniform}
 \int_{\mathbb R}|\widehat w(\sigma+i\tau)|
 \left|\sum_{\ell=1}^H b_\ell M_\ell(\sigma+i\tau)\right|d\tau
       \ll_\sigma (T+H)\log(2H).
\end{equation}

Here are details of the gamma estimate and its summation, to make
the frequency uniformity explicit. Write
$p_{j,\ell}=a_{j,\ell}D_{s_j}(\kappa_\ell)$ as in
\eqref{A:Reznikov-mean}. Put
$a=3/4-\sigma/2$ and $b=1/2-\sigma/2$, so $1/2<a<3/4$.
For $\tau,t,\kappa\geq0$ Stirling's formula gives
\begin{align}\label{A:gamma-full-bound}
 |G(\sigma+i\tau,t,\kappa)D_{1/2+it}(\kappa)|
 &\ll_\sigma U^{1/2}
  \{(1+\tau+t)(1+|\tau-t|)\}^{-a}\\
 &\quad\cdot\{(1+\tau+\kappa)(1+|\tau-\kappa|)\}^{b}
 \\ &\quad\cdot\{(1+t+\kappa)(1+|t-\kappa|)\}^{-1/4}
 e^{-c\operatorname{dist}(t,[\min(\tau,\kappa),\max(\tau,\kappa)])}.
 \nonumber
\end{align}
One may take $c=\pi/2$. The exponential follows exactly from
\[
 t-\max(t,\kappa)-\max(t,\tau)+\max(\tau,\kappa)
 =-\operatorname{dist}(t,[\min(\tau,\kappa),\max(\tau,\kappa)]).
\]
The same bound holds with $|\tau|$, $|\kappa|$ for arbitrary signs.
The fixed real parts of the numerator gamma functions are
$\sigma/2-1/4$, which are negative but not nonpositive integers;
thus the estimate is bounded at $\tau=\pm t$.

In a dyadic spectral band $J\leq1+t_j<2J$, the elementary
estimates and Reznikov's bound used in Statement A.4.3 give
\begin{equation}\label{A:three-basic-sums}
 \sum_j|p_j'|^2\ll J,\quad
 \sum_j|p_j'|\ll J^{3/2},\quad
 \sum_j|p_j'p_{j,\ell}|\ll J,
 \quad\sum_{|\ell|\leq K}|a_{j,\ell}|^2\ll\max(K,J).
\end{equation}
The middle bound uses the elementary $O(J^2)$ eigenvalue count.
Changing between $\ell$ and $\kappa_\ell$ changes dyadic bounds
only by fixed constants.

First suppose $\kappa_\ell\leq t_j/2$. Then
$|D_{s_j}(\kappa_\ell)|\asymp J^{-1/2}$.
For $\tau\geq t_j/2$, \eqref{A:gamma-full-bound} implies
\[
 |G(\sigma+i\tau,t_j,\kappa_\ell)|
 \ll_\sigma
 \left(\frac{1+\tau}{1+|\tau-t_j|}\right)^a
 e^{-c(t_j-\tau)_+}.
\]
For $\tau<t_j/2$ there is instead exponential decay in $J$,
with an immaterial polynomial factor. Since $a<1$, integration
over a dyadic $\tau$ block of size $U\gg J$, with the factor
$U^{-1}(1+U/T)^{-N}$, is
$O((1+U/T)^{-N})$; the same assertion holds for $U\asymp J$
because the integrable singular factor is bounded by
$(J/(1+|\tau-t_j|))^a$.
Summing \eqref{A:three-basic-sums} over $\ell$ costs $\log(2H)$.
Summing over $U\gg J$ and then dyadic $J$ costs $O(T)$:
for $J\leq T$ the contribution is
$O(J\log(2T/J))$, whose dyadic sum is $O(T)$, and the remaining
bands decay faster than any power of $J/T$.

Second suppose $\kappa_\ell\geq2t_j$, in an angular block
$K\leq\kappa_\ell<2K$ with $K\gg J$.
Up to $\tau\leq\max(2t_j,\kappa_\ell/2)$,
\eqref{A:gamma-full-bound} is bounded by
\[
 K^{1/2-\sigma}U^{1/2}J^{-a}(1+|\tau-t_j|)^{-a}
 e^{-c(\tau-t_j)_+}.
\]
In the remaining range, the exponential factor is at most
$e^{-c'K}$. For $\tau\gg K$ the polynomial factor is bounded
independently of $\tau$ (the total $\tau$ power in
\eqref{A:gamma-full-bound} is zero). For $\tau\asymp K$ it is
at most a fixed power of $K$. After the coefficient bounds below,
the sum of all these tails is therefore $O(\log(2T))$:
the integral of $U^{-1}(1+U/T)^{-N}$ above a fixed positive
lower limit is $O(\log(2T))$, and the exponential in $K$
absorbs the polynomial factors. This accounts also for
$\tau\geq\kappa_\ell$, where the exponential is
$e^{-c(\kappa_\ell-t_j)}$, rather than
$e^{-c(\tau-t_j)}$.
Cauchy--Schwarz in $\ell$, using
$\sum|a_{j,\ell}|^2\ll K$ and
$\sum_{\ell\asymp K}|b_\ell|^2\ll K^{-1}$, introduces no
additional factor. The elementary beta-integral bound
\[
 \int_0^\infty U^{-1/2}(1+|\tau-t_j|)^{-a}
 e^{-c(\tau-t_j)_+}(1+U/T)^{-N}\,d\tau
 \ll_\sigma J^{-a}\min(J,T)^{1/2}
\]
therefore bounds the spectral band by
\[
 K^{1/2-\sigma}J^{3/2-2a}\min(J,T)^{1/2}
 =K^{1/2-\sigma}J^\sigma\min(J,T)^{1/2}.
\]
Its dyadic sum over $J\ll K$ is $O_\sigma(K)$ for $K\leq T$
and $O_\sigma(\sqrt{KT})$ for $K>T$. Summing $K\ll H$ is
$O_\sigma(T+H)$.

Finally suppose $t_j/2<\kappa_\ell<2t_j$, so $J\ll H$.
Cauchy--Schwarz in $\ell$ gives
\[
 \left|\sum b_\ell a_{j,\ell}
       G(\sigma+i\tau,t_j,\kappa_\ell)D_{s_j}(\kappa_\ell)\right|
 \ll J^{-1/2}
       \left(\sum_{\kappa_\ell\asymp J}|GD|^2\right)^{1/2}.
\]
Squaring \eqref{A:gamma-full-bound} and using
$\sum_{\kappa_\ell\asymp J}(1+|t_j-\kappa_\ell|)^{-1/2}
\ll J^{1/2}$ gives, in the three respective $\tau$ ranges,
\[
 \sum_{\kappa_\ell\asymp J}|GD|^2\ll_\sigma
 \begin{cases}
 U/J,&\tau\leq t_j/2,\\
 J^{2b+1/2}(1+|\tau-t_j|)^{-2a},&t_j/2<\tau<2t_j,\\
 1,&\tau\geq2t_j.
 \end{cases}
\]
After multiplication by $|\widehat w|$, integration in $\tau$,
and summation of $|p_j'|$ over the spectral band, these give
respectively
\[
 O_\sigma(\sqrt{J\min(J,T)}),\qquad
 O_\sigma(J(1+J/T)^{-N}),\qquad
 O_\sigma\bigl(J\log(2+T/J)\min(1,(T/J)^N)\bigr).
\]
For the middle range the exponent is exactly
$b-1/4-a=-1/2$ before the final $J^{3/2}$ sum.
Their dyadic sum over $J\ll H$ is $O_\sigma(T+H)$.
These three cases prove \eqref{A:contour-uniform}. The finitely
many $0<t_j<1$ are included in the band $J=1$; the strict spectral
gap ensures all constants there are finite.

\emph{4. Shifting the contour and bounding the residues.}
For $Y\geq2$ write
\[
 W_\ell(Y)=\sum_{z\geq1}\sum_{N\mathfrak a=m_z}
 e^{2\pi i\ell\theta(\mathfrak a)}w(A_z/Y).
\]
Mellin inversion of \eqref{A:Mellin-geometry} on a line
$\operatorname{Re}s>1$ gives $4W_\ell$ as the inverse Mellin
transform of $M_\ell-2(-1)^\ell-4H_\ell$.
For $H_\ell$ shift to $\operatorname{Re}s=\sigma$ without
crossing a pole. Equations \eqref{A:H-uniform} and
\eqref{A:w-Mellin-bound} show, uniformly in the coefficients,
\begin{equation}\label{A:correction-inversion}
 \left|\sum_{\ell=1}^Hb_\ell\frac1{2\pi i}
 \int_{(\sigma)}\widehat w(s)Y^sH_\ell(s)\,ds\right|
 \ll_\sigma Y^\sigma
     T^{1/2}(T+H)^{1-\sigma}\log(2H).
\end{equation}
Indeed the remaining integral is bounded by
$\int_0^\infty U^{-1/2}(U+H)^{1-\sigma}(1+U/T)^{-N}\,d\tau$.
The exceptional constant $2(-1)^\ell$ contributes
$2(-1)^\ell w(1/Y)$ and hence $O(\log(2H))$ in the angular sum.

Shift the $M_\ell$ contour to the same line. Its residue at $s=s_j$ is
\[
 p_j'p_{j,\ell}\,r_j\frac{D_{s_j}(0)}{D_{s_j}(\kappa_\ell)},\qquad
 r_j=\frac{2\Gamma((s_j+1)/2)\Gamma(s_j-1/2)}{\Gamma(s_j/2)^3},
 \qquad |r_j|\ll(1+t_j)^{1/2}.
\]
The residue at $1-s_j$ has the conjugate gamma factors and obeys
the same estimate. Statement A.4.3 and
\eqref{A:w-Mellin-bound} therefore bound all residues, after summing
the angular coefficients, by
\begin{equation}\label{A:residues-inversion}
 \ll Y^{1/2}\log(2H)
      \sum_{J\text{ dyadic}}J^{1/2}(1+J/T)^{-N}
 \ll Y^{1/2}T^{1/2}\log(2H).
\end{equation}
The lower-half-plane residues use the same proof with conjugated
gamma ratios. No positivity or common phase of the periods is used.
The remaining contour costs
$O_\sigma(Y^\sigma(T+H)\log(2H))$ by
\eqref{A:contour-uniform}.

For clarity, all contour shifts can first be made at finite heights.
The eigenvalue bound $N(V)\ll V^2$ permits heights tending to infinity
whose distance from every $\pm t_j$ is at least a fixed multiple of
$V^{-3}$: the excluded intervals in $[V,V+1]$ have total length
$O(V^{-1})$. On these horizontal segments the gamma formula and
the spectral coefficient bounds give polynomial growth for each
fixed $H$. Taking $N$ arbitrarily large in
\eqref{A:w-Mellin-bound} makes their integrals tend to zero.
Absolute convergence of the residue series follows from
\eqref{A:residues-inversion}. The estimates on the resulting vertical
line are uniform in $H,T$, so this limiting argument makes no
uncontrolled exchange of a growing angular cutoff with a contour limit.

Combining these estimates proves the uniform smoothed inequality
\begin{equation}\label{A:smoothed-angular-sum}
 \left|\sum_{\ell=1}^H b_\ell W_\ell(Y)\right|
 \ll_\sigma\log(2H)\left[
 Y^{1/2}T^{1/2}+Y^\sigma(T+H)
 +Y^\sigma T^{1/2}(T+H)^{1-\sigma}+1\right].
\end{equation}
Because the coefficients $b_\ell$ are arbitrary within
$|b_\ell|\leq1/\ell$, choose their phases to align $W_\ell(Y)$.
Then the left side is exactly
$\sum_{\ell\leq H}|W_\ell(Y)|/\ell$.
This averaged angular estimate is sufficient; it avoids asserting
the stronger, individual-mode estimate \eqref{A:needed-twist}.

\emph{5. Returning to both sharp boundaries.}
The Erd\H{o}s--Tur\'an inequality used in Statement A.4.2 applies
also to a finite list with nonnegative real weights. For rational
weights this follows by replication after clearing denominators;
arbitrary weights follow by approximation and continuity. Hence the
weighted arc discrepancy is at most
\[
 O\left(\frac{W_0(Y)}H+
       \sum_{\ell\leq H}\frac{|W_\ell(Y)|}{\ell}+\log(2Y)\right).
\]
The endpoint term is $O(\log(2Y))$ by the Pell calculation already
proved; the other arc endpoint is empty.
Since Statement A.3.1 gives
$\#\{(z,\mathfrak a):A_z\leq V\}=(c/3)V+O(V^{2/3})$,
partial summation gives
\[
 W_0(Y)=\frac c3Y\widehat w(1)+O(Y^{2/3}),
 \qquad c=\frac{9L_0^2}{2\pi^2},
\]
uniformly for the monotone smooth cutoffs under discussion, because
their total variation is one.

Choose two such cutoffs $w^-,w^+$ with
$0\leq w^-\leq\mathbf1_{[0,1]}\leq w^+\leq1$,
transition intervals respectively $[1-1/T,1]$ and $[1,1+1/T]$.
They bracket the sharp radial count with nonnegative weights and
satisfy $\widehat w^\pm(1)=1+O(1/T)$.
Put $T=H=\lfloor Y^{1/3}\rfloor$ in
\eqref{A:smoothed-angular-sum}. The three nonconstant error terms
are respectively
\[
 O(Y^{2/3}\log Y),\quad
 O_\sigma(Y^{1/3+\sigma}\log Y),\quad
 O_\sigma(Y^{1/2+2\sigma/3}\log Y).
\]
The angular and radial main-term bracketing errors are
$O(Y/H+Y/T)=O(Y^{2/3})$.
Take any fixed sufficiently small $\sigma>0$, and absorb logarithms
into $Y^\epsilon$. Finally put $Y=(9X^2-1)^{1/2}$.
The ideal-arc main term is $(\lambda_0c/3)Y=C_0X+O(X^{-1})$.
Statement A.1.3 changes this by only $O(\log(2X))$, proving
\eqref{A:sharp-two-thirds}.
\end{proof}


\paragraph{Statement A.4.5 (Proved endpoint-strip rigidity).}
Put $b=b_0$, $A=A_z=\sqrt{9z^2-1}$ and
$\delta=\operatorname{arsinh}(1/A)$. For an integral solution
$s^2-3t^2=2A^2$, $s>0$, write
$s=\sqrt2A\cosh u$, $t=\sqrt{2/3}A\sinh u$.
After choosing the lift of $u$ nearest the indicated endpoint,
\[
 |u|\leq\delta\ \Longrightarrow\ t=0,
 \qquad |u-b|\leq\delta\ \Longrightarrow\
 \bigl(4s-17z=0\text{ or }(s,t,z)=(4,0,1)\bigr).
\]
Consequently only $O(\log(2X))$ reduced ideals with $z\leq X$
can lie within these two endpoint strips. These are exactly the
possible endpoint-crossing terms in the spherical correction of
Statement A.4.4; this assertion concerns the geometric support of
that correction, not its complete amplitude.
\begin{proof}
The first implication follows from
$|t|\leq\sqrt{2/3}A\sinh\delta=\sqrt{2/3}<1$.
For the second use
\[
 \cosh b=\frac{17}{12\sqrt2},\qquad
 \sinh b=\frac1{12\sqrt2},\qquad
 \cosh\delta=\frac{3z}{A}.
\]
It follows that
\[
 \sqrt2A\cosh(b\pm\delta)=\frac{17z}{4}\pm\frac1{12}.
\]
Since $b>0$ and $\cosh$ is increasing on $[0,\infty)$, this
bounds $s$ when $\delta<b$, since the entire strip is positive.
In the remaining case $\delta\geq b$, we have $A\leq
1/\sinh b=12\sqrt2$, hence $z\leq5$. Directly for these five
heights, the only integral solutions in the second strip are
$(4,0,1)$ and $(17,1,4)$; equivalently one can note that the negative part lies
in the first strip and has $t=0$, and test $s^2=18z^2-2$ for
$z=1,2,3,4,5$. The only $t=0$ point in this finite range is $(4,0,1)$.
Outside this one stated exception, $|4s-17z|\leq1/3$, and integrality forces
$4s=17z$. The latter equation gives $z=4a$, $s=17a$ and
$a^2-3t^2=-2$, precisely the $n=0$ Pell family already counted.
The first equation gives $s^2-18z^2=-2$, another Pell equation;
its positive solutions form finitely many unit orbits, hence have
$O(\log(2X))$ members up to height $X$. Alternatively, ordinary
reduction by the fixed unit $17+12\sqrt2$ proves this bound directly.

Finally, in \eqref{A:F-spherical-integral} the multiplier of
$e^{i\kappa u}$ is an integral of
$e^{-i\kappa\log B(x)/2}$, with
$B(x)=\cosh(2\delta)+x\sinh(2\delta)$ and $-1\leq x\leq1$.
Since $e^{-2\delta}\leq B(x)\leq e^{2\delta}$, it only shifts
$u$ by a number in $[-\delta,\delta]$. An interval indicator
can therefore change only in the strips just considered. There
is also the scalar factor $B(x)^{-s/2}$: away from the strips
it gives $F_{s,0}(A)$, rather than $1$. No disappearance of that
amplitude correction is claimed.
\end{proof}

\paragraph{Statement A.4.6 (Proved existence of the sharp spectral arc functional).}
For each fixed nonconstant eigenfunction $u_j$, the limit
\[
 Q_{j,I}=\lim_{H\to\infty}
 \sum_{|\ell|\leq H}\widehat{\mathbf1_I}(\ell)
 \frac{D_{s_j}(0)}{D_{s_j}(\kappa_\ell)}p_{j,\ell}
\]
exists. Here $I=[0,b_0/L_0]$, the Fourier coefficient is with
respect to $\theta=u/L_0$, and the periods are unnormalized as
in Statement A.4.4. The limit is a well-defined interval value
of a distribution whose local primitive is H\"older continuous
with exponent $1/2$. This proves existence, not nonvanishing,
of the actual arc coefficient.
\begin{proof}
Conjugate the lift of the source axis to $\{ie^u:u\in\mathbb R\}$
in the upper half-plane. The Poisson--Helgason boundary
representation writes the lifted eigenfunction as
\[
 u_j(z)=\int_{\mathbb R}
       \left(\frac{\operatorname{Im}z}{|z-x|^2}\right)^{s_j}
       T_j(dx).
\]
This is a distributional integral, with the usual additional
chart at infinity. The precise theorem used is the boundary-value
representation in Anantharaman--Zelditch
\cite[Theorem 3.2]{AnantharamanZelditch2012}, together with their
Theorem 7.1 (Otal's boundary regularity theorem): a bounded
hyperbolic Laplace eigenfunction with bounded gradient and
parameter $s$, $\operatorname{Re}s\geq0$, has boundary values
which are the distributional derivative of a
$\operatorname{Re}s$-H\"older function. Our eigenfunction and
its gradient are bounded because the quotient is compact;
Statement A.2.5 gives $\operatorname{Re}s_j=1/2$. The source is
written for compact torsion-free quotients, so apply it to the
fixed cover of Statement A.3.1. The lifted eigenfunction is
unchanged as a function on the universal cover.

For $x=\pm e^v$, define the sum of two boundary distributions
\[
 \beta_j(dv)=(2e^v)^{-s_j}
       \bigl(T_j(d(e^v))+T_j(d(-e^v))\bigr),
\]
where $T_j(d(-e^v))$ denotes pushforward under
$x\mapsto\log|x|$ on the negative half-line, without an
orientation sign. Smooth changes of variable and multiplication
by a smooth function preserve the property of being the
derivative of a locally $1/2$-H\"older function. To see the
multiplication assertion explicitly, if $T=F'$ then
$aT=(aF)'-a'F$, and the second term has a continuously
differentiable primitive. The same formula after a smooth
coordinate change proves the other assertion.

The dilation $x\mapsto e^{L_0}x$ belongs to the source
stabilizer. Uniqueness of the Poisson boundary representation
implies $T_j(e^{L_0}dx)=e^{L_0s_j}T_j(dx)$, so $\beta_j$ is
$L_0$-periodic. Since
\[
 \left(\frac{e^u}{e^{2v}+e^{2u}}\right)^{s_j}
     =(2e^v)^{-s_j}\cosh(u-v)^{-s_j},
\]
its convolution with $\cosh(u)^{-s_j}$ is the restriction
$u_j(ie^u)$. The convolution is legitimate for the periodic
distribution because this kernel and each of its derivatives
decay exponentially. Taking its Fourier coefficients gives
\[
 p_{j,\ell}=K_{s_j}(\kappa_\ell)\,\beta_{j,\ell},
 \qquad K_s(\kappa)=2^{s-1}D_s(\kappa),\qquad
 \beta_{j,\ell}=\int_0^{L_0}e^{i\kappa_\ell v}\beta_j(dv).
\]
Therefore the displayed partial sum defining $Q_{j,I}$ equals
$K_{s_j}(0)$ times the pairing of $\beta_j$ with the symmetric
Fourier partial sum of $\mathbf1_{[0,b_0]}$.

Let $F_j$ be a continuous local primitive of $\beta_j$;
subtracting its affine mean gives a periodic $1/2$-H\"older
function. Its Fourier partial sums converge uniformly. For
completeness, convolving a H\"older function with the Fej\'er
kernel gives uniform error $O_j(H^{-1/2})$ by splitting the
kernel integral at $1/H$; composing the approximation with the
Dirichlet projection of degree $H$ and using its elementary
$L^1$ bound $O(\log(2H))$ gives an error
$O_j(H^{-1/2}\log(2H))$ for the ordinary partial sum (use a
Fej\'er polynomial of degree at most $H$). Integration by parts
then identifies the limit with
\[
 Q_{j,I}=K_{s_j}(0)\bigl(F_j(b_0)-F_j(0)\bigr).
\]
This value is independent of the primitive's additive constant,
and the convergence error is $O_j(H^{-1/2}\log(2H))$.
The implied constant here depends on the fixed eigenfunction;
no exchange with an unbounded spectral sum follows from this
fixed-eigenfunction assertion alone.
\end{proof}

\paragraph{Statement A.4.7 (Proved: exact smooth model for this particular sharp arc).}
This statement concerns the actual arc, and uses its arithmetic endpoints.
Put
\[
 a_*=17,\qquad B_*=12\sqrt2,\qquad
 R_*=\operatorname{arcosh}(17)=2\operatorname{arcosh}(3),
 \qquad b=\log\frac3{2\sqrt2}.
\]
Thus $a_*^2-B_*^2=1$ and $\coth R_*=\cosh b$.
For $\Re s>0$, define
\[
 W_s(u)=(a_*+B_*\cosh u)^{-s},\qquad
 H_s(\kappa)=\int_{\mathbb R}W_s(u)e^{i\kappa u}\,du,
 \qquad m_s(\kappa)=H_s(\kappa)/K_{s,\kappa}.
\]
Choose a real eigenbasis also diagonalizing the reflection involution.
Every mode that contributes to $p'_j$ is reflection-even. For those modes,
\begin{align}\label{A:U-smooth-boundary}
 \beta_j([0,b])
 &=\frac12\int_{\mathbb R}W_{s_j}(u)\,\beta_j(du),\\
 Q_{j,I}
 &=\frac{K_{s_j,0}}{2L_0}
       \sum_{\ell\in\mathbb Z}p_{j,\ell}\,
                    m_{s_j}(\kappa_\ell).
 \label{A:U-smooth-Q}
\end{align}
The first integral is the distributional pairing with a smooth,
exponentially decreasing function; the second series converges absolutely.
In particular no truncation at the sharp endpoints occurs in
\eqref{A:U-smooth-Q}.

The multiplier has the exact formula
\begin{align}\label{A:U-Legendre-multiplier}
 H_s(\kappa)
 &=\sqrt{\frac{2\pi}{B_*}}\,
   \frac{\Gamma(s+i\kappa)\Gamma(s-i\kappa)}{\Gamma(s)}
   P_{-1/2+i\kappa}^{\,1/2-s}(a_*/B_*),\\
 m_s(\kappa)
 &=\frac{2^{s-1/2}}{\sqrt{\pi B_*}}\,
   \Gamma\!\left(\frac{s+1+i\kappa}{2}\right)
   \Gamma\!\left(\frac{s+1-i\kappa}{2}\right)
   P_{-1/2+i\kappa}^{\,1/2-s}(a_*/B_*).
 \nonumber
\end{align}
For $s=1/2+it$, $t\geq0$, and $|\kappa|\geq t$,
\begin{equation}\label{A:U-forbidden-exponential}
 |m_s(\kappa)|\ll
 (1+t+|\kappa|)^{1/2}
       e^{-\pi(|\kappa|-t)/4}.
\end{equation}

For each fixed $\delta>0$ there is also a uniform stationary-phase
description on $|\rho|=|\kappa/t|\leq1-\delta$ as $t\to\infty$.
Set
\[
 q(\rho)=\sqrt{B_*^2+\rho^2},\qquad
 \xi(\rho)=
 \log\frac{q(\rho)+a_*\rho}{B_*\sqrt{1-\rho^2}},
\]
and
\[
 S(\rho)=\log(a_*+q(\rho))
       -\tfrac12\log(1-\rho^2)-\rho\,\xi(\rho).
\]
Then
\begin{equation}\label{A:U-stationary-multiplier}
 m_{1/2+it}(\kappa)
 =q(\rho)^{-1/2}(1-\rho^2)^{1/4}
       e^{-itS(\rho)}\{1+O_\delta(t^{-1})\}.
\end{equation}
Its canonical time and tangential displacement obey the exact identities
\begin{equation}\label{A:U-canonical-time}
 S'(\rho)=-\xi(\rho),\qquad
 S(\rho)-\rho S'(\rho)
   =\operatorname{arsinh}\!\bigl(B_*\cosh\xi(\rho)\bigr)
   \geq R_*.
\end{equation}
These formulas and the exponential forbidden-region estimate do not by
themselves assert nonvanishing of an eigenspace contraction.

\begin{proof}
All reflections represent the same nontrivial coset of
$W/\Gamma^+$. Thus an odd eigenfunction vanishes pointwise on every
reflection mirror, in particular on both target geodesics.
Its two target periods vanish. We may therefore restrict to
$W$-even eigenfunctions.

In the upper half-plane model with source axis $i e^u$, the reflections
in $z=0$ and $t=0$ act on its real boundary as $x\mapsto-x$ and
$x\mapsto1/x$, respectively. Consequently the boundary distribution
$T_j$ in Statement A.4.6 satisfies $T_j(-dx)=T_j(dx)$, and
$\beta_j$ is even in $u$. The element
\[
 U=S_4S_1,\qquad U(x)=\frac{3x+2\sqrt2}{2\sqrt2 x+3}
\]
belongs to $\Gamma^+$, fixes $x=1$, and maps $\infty$ to $e^b$.
The boundary transformation law, obtained directly from the Poisson
kernel and uniqueness of its distributional boundary value, is
$T_j(U\,dx)=|U'(x)|^{s_j}T_j(dx)$.
On the positive boundary, the two equal signs of $x$ give
$\beta_j(du)=2(2e^u)^{-s_j}T_j(d(e^u))$.
Changing variables $x=U(y)$ on the interval $[1,e^b]$ therefore gives
\[
 \beta_j([0,b])
 =\int_0^\infty
 \left(\frac{e^u}
 {(3e^u+2\sqrt2)(2\sqrt2e^u+3)}\right)^{s_j}
 \beta_j(du)
 =\int_0^\infty W_{s_j}(u)\,\beta_j(du).
\]
Evenness gives \eqref{A:U-smooth-boundary}. The change of variables
is first made away from the endpoints. The boundary primitive is
$1/2$-H\"older, so its interval masses are well-defined; the tail
of the transformed pairing tends to zero because $W_s$ and all its
derivatives decrease exponentially for fixed $\Re s>0$.
This justifies passage to the endpoints.
Periodize $W_s$ with period $L_0$ and use
$p_{j,\ell}=K_{s_j,\kappa_\ell}\beta_{j,\ell}$.
The smooth periodic kernel has rapidly decreasing Fourier coefficients,
which proves \eqref{A:U-smooth-Q}; the stronger convergence assertion
also follows from \eqref{A:U-forbidden-exponential}.

For the special-function identity, take
$x=a_*/B_*>1$, $\nu=-1/2+i\kappa$, and $\mu=s-1/2$
in DLMF \cite[equation (14.12.4)]{DLMF}. Its conditions
$\Re(\mu-\nu)>0$ and $\nu+\mu\notin\{-1,-2,\ldots\}$
hold for $\Re s>0$ and real $\kappa$.
Since $x^2-1=B_*^{-2}$, its integral is exactly one half
of $B_*^sH_s(\kappa)$. This gives the first formula.
The gamma duplication formula gives the second.
For $\kappa=0$, equation (14.12.6), with the conversion
from Olver's $\mathbf Q$ to the standard $Q$, gives
$H_s(0)=2Q_{s-1}(17)$. For example
$H_1(0)=2Q_0(17)=2b$, verifying the constant-mode normalization.

For the uniform bound, the analytic continuation of $W_s(u)$
has no singularity in $|\Im u|<\pi$: its first zeros are
at $\pm b\pm i\pi$.
Shift the contour to $\Im u=3\pi/4$ when $\kappa\geq0$,
and to $\Im u=-3\pi/4$ when $\kappa<0$.
For $v=3\pi/4$ and real $u$,
\[
 |a_*+B_*\cosh(u+iv)|^2
   =288\cosh^2u-408\cosh u+145
   \geq25\cosh^2u.
\]
Also $|\arg(a_*+B_*\cosh(u+iv))|\leq v$.
For $u\geq0$ the imaginary part is nonnegative, and
the imaginary part after multiplication by $e^{-iv}$ equals
$-\sin v\,(a_*+B_*e^{-u}\cos v)<0$;
the other signs follow by conjugation. Hence
\[
 |H_{1/2+it}(\kappa)|
 \leq 5^{-1/2}K_{1/2,0}
       e^{-3\pi(|\kappa|-t)/4}.
\]
The vertical connecting segments vanish by exponential decay in
$|\Re u|$. Uniform two-sided Stirling estimates give
\[
 |K_{1/2+it,\kappa}|
 \asymp
 [(1+|t+\kappa|)(1+|t-\kappa|)]^{-1/4}
 e^{-\frac\pi4(|t+\kappa|+|t-\kappa|-2t)}.
\]
Combining the two estimates proves
\eqref{A:U-forbidden-exponential}.

Finally apply stationary phase separately to $H$ and $K$.
The phases are $\rho u-\log(a_*+B_*\cosh u)$ and
$\rho u-\log\cosh u$.
Their unique stationary points are
\[
 u_W=\log\frac{q(\rho)+a_*\rho}{B_*(1-\rho)},
 \qquad u_K=\operatorname{artanh}\rho .
\]
On $|\rho|\leq1-\delta$ they stay in a fixed compact set and
are uniformly nondegenerate. Outside a larger fixed compact set
the phase derivatives are uniformly separated from zero;
integration by parts treats these exponentially decreasing tails.
At the stationary points, the amplitude divided by the square root
of the absolute second phase derivative equals $q(\rho)^{-1/2}$
for $H$, and $(1-\rho^2)^{-1/4}$ for $K$.
The two stationary-phase factors have the same sign and cancel
in the quotient. This proves \eqref{A:U-stationary-multiplier}.
The formulas for $S'$ and the canonical time follow by differentiation.
Explicitly
$\cosh\xi=q/(B_*\sqrt{1-\rho^2})$ and
$\sinh\xi=a_*\rho/(B_*\sqrt{1-\rho^2})$; substitution gives
\eqref{A:U-canonical-time}.
\end{proof}

\paragraph{Scope of this additional identity.}
The axis of $U$ is the mirror $t=0$, whose normal has norm $-3$.
It is not either norm-$-2$ target.
The two roots of the smooth kernel, $-e^{\pm b}$, are endpoints
of $U^{-1}g_0$, a translate of the source geodesic.
The positive boundary component used in the integral is not a
full period fundamental interval for that axis. Thus the identity
does not turn $Q_{j,I}$ into a positive scalar multiple of $p'_j$.
The nonzero complete-eigenspace coefficient required in
Statement A.5.6 is proved separately in Statement A.5.8, using
the outgoing physical kernel and the complete endpoint comparison.


\paragraph{Statement A.4.8 (Proved: exact outgoing physical kernel, including glancing).}
Retain $a_*,B_*,R_*$ from Statement A.4.7 and put
\[
 \tau(u)=\operatorname{arsinh}(B_*\cosh u),\qquad
 f_s(r)=2^{-s}(\cosh r)^{-s}
 {}_2F_1\!\left(\frac s2,\frac s2;s+\frac12;\operatorname{sech}^2r\right).
\]
For $\Re s>0$ the inverse Fourier kernel of $m_s$ is
\begin{equation}\label{A:exact-physical-kernel}
 k_s(u)=-\frac{\Gamma(s)}{\sqrt\pi\,\Gamma(s+1/2)}
          \frac{a_*}{\cosh\tau(u)}f_s'(\tau(u)).
\end{equation}
Here the convention is
$m_s(\kappa)=\int_{\mathbb R}k_s(u)e^{i\kappa u}\,du$.
For every contributing reflection-even eigenfunction the exact sharp
coefficient is therefore
\begin{equation}\label{A:Q-physical}
 Q_{j,I}=\int_{\mathbb R}u_j(i e^u)\,q_{t_j}(u)\,du,
 \qquad q_t(u)=\frac{K_{1/2+it,0}}2 k_{1/2+it}(u).
\end{equation}
The integral is absolutely convergent. As $t\to+\infty$,
\begin{equation}\label{A:physical-principal-symbol}
 q_t(u)=
 \frac{a_*}{2\cosh^{3/2}\tau(u)}
       e^{-it\tau(u)}\{1+O(t^{-1})\},
 \qquad u\in\mathbb R,
\end{equation}
uniformly over the entire real line.
The principal amplitude is positive and decreases like
$e^{-3|u|/2}$. Thus the physical formula remains uniform where
$|\kappa|/t$ tends to one in the Fourier description.

\begin{proof}
The hypergeometric equation gives directly
\[
 f_s''+\tanh r\,f_s'+s(1-s)f_s=0,\qquad
 f_s(r)\sim e^{-sr}\quad(r\to+\infty).
\]
We prove the Fourier identity first for real $s>1$ and then use
analytic continuation. In this proof temporarily allow
$a_*=\cosh R$, $B_*=\sinh R$ with arbitrary $R>0$.
The elementary function $W_s(R,u)$ satisfies
\[
 \left(\partial_R^2+\coth R\,\partial_R+s(1-s)
              -\frac{\partial_u^2}{\sinh^2R}\right)W_s(R,u)=0.
\]
To check the candidate kernel, set $y=\sinh r$ and write $f_s(r)=F_s(y)$.
Then
\[
 (1+y^2)F_s''+2yF_s'+s(1-s)F_s=0.
\]
Consequently $g=F_s'$ satisfies
\[
 (1+y^2)g''+4yg'+\{s(1-s)+2\}g=0.
\]
A direct differentiation shows that
$\cosh R\,g(\sinh R\cosh u)$ satisfies the same two-variable
equation as $W_s$. Its scalar multiple in
\eqref{A:exact-physical-kernel} therefore does so as well.
Fourier transformation in $u$ makes both transforms solutions of
\[
 h''+\coth R\,h'
       +\left(s(1-s)+\frac{\kappa^2}{\sinh^2R}\right)h=0.
\]

We compare their outgoing coefficients. Uniformly under the integrable
majorants needed for Fourier transformation,
\[
 W_s(R,u)\sim e^{-sR}\cosh^{-2s}(u/2),\qquad
 k_s(R,u)\sim
 \frac{s\Gamma(s)}{\sqrt\pi\,\Gamma(s+1/2)}
 e^{-sR}\cosh^{-s-1}u.
\]
The duplication formula gives the exact identity
\[
 \frac{K_{s,\kappa}K_{s+1,\kappa}}
            {2K_{2s,2\kappa}}
   =\frac{\sqrt\pi\,\Gamma(s+1/2)}{s\Gamma(s)}.
\]
It follows that the Fourier transform of $k_s$ and $H_s/K_s$
have the same leading $e^{-sR}$ coefficient.
For $s>1$, the other solution grows like $e^{(s-1)R}$;
the outgoing solution is unique.
Thus their transforms coincide.
All defining integrals are locally holomorphic for $\Re s>0$,
and $K_{s,\kappa}$ has neither zeros nor poles there.
The identity theorem proves \eqref{A:exact-physical-kernel}
throughout that half-plane. Statement A.4.7 and Fourier
periodization now give \eqref{A:Q-physical}.

For the uniform high-frequency assertion, write
$v_t(r)=\sqrt{\cosh r}\,f_{1/2+it}(r)$.
The equation becomes the elementary one-dimensional scattering equation
\[
 v_t''+\left(t^2-\frac14\operatorname{sech}^2r\right)v_t=0,
 \qquad v_t(r)\sim 2^{-1/2}e^{-itr}.
\]
On $r\geq R_*>0$, its outgoing Volterra equation has an
integrable smooth potential $V(r)=\operatorname{sech}^2r/4$.
Successive approximation, using the Green kernel
$t^{-1}\sin(t(r-v))$ and
$\int_r^\infty V(v)\,dv\ll e^{-2r}$, yields
\[
 v_t(r)=2^{-1/2}e^{-itr}
       \{1+O(t^{-1}e^{-2r})\},\qquad
 v_t'(r)=-it\,2^{-1/2}e^{-itr}+O(e^{-2r}).
\]
The constants are uniform for $t\geq1$ and $r\geq R_*$.
Therefore
\[
 f_{1/2+it}'(r)=
 -\frac{it}{\sqrt{2\cosh r}}e^{-itr}\{1+O(t^{-1})\}.
\]
If further finite-order symbol expansions are required, they follow
from the same Volterra equation: after removing $e^{-itr}$,
insert $\sum_{n=0}^N(2it)^{-n}a_n(r)$, with $a_0=1$ and
$a_{n+1}'=a_n''-Va_n$, $a_{n+1}(\infty)=0$.
The residual is $O(t^{-N}e^{-2r})$, and the outgoing Green kernel
gains a further factor $t^{-1}$. All coefficients and their fixed-order
$r$ derivatives are bounded, and the nonconstant coefficients decrease
exponentially. This also verifies that the preceding uniform
estimate is compatible with finite-order stationary-phase calculations.

Finally, Stirling's formula or direct stationary phase in $K$ gives
\[
 \frac{\Gamma(1/2+it)}{\Gamma(1+it)}
       =t^{-1/2}e^{-i\pi/4}\{1+O(t^{-1})\},\qquad
 K_{1/2+it,0}=
       \sqrt{\frac{2\pi}{t}}e^{-i\pi/4}\{1+O(t^{-1})\}.
\]
The factor $it$ from $f_s'$ cancels both phases.
Substitution in \eqref{A:exact-physical-kernel} proves
\eqref{A:physical-principal-symbol}. Since
$\cosh\tau(u)\asymp e^{|u|}$, it also proves the stated decay.
\end{proof}

\paragraph{Connection with the nonvanishing argument.}
The exact physical kernel makes the two pieces of a relative-wave
argument concrete. On bounded $u$ intervals, the phase is the signed
difference of distances to the target geodesic and to $U^{-1}g_0$;
its critical points represent common perpendiculars. At infinity,
these differences approach endpoint Busemann constants.
Those constants can accumulate, so their contribution requires a
quantitative summation before the shortest-perpendicular term
proves nonvanishing. Statement A.5.8 supplies that comparison:
the entire endpoint contribution is the Fourier transform of a
bounded-variation density with finite total variation, and is
of lower order than the positive stationary term.


\subsection{Exact census and square-root-scale evidence}
\paragraph{Statement A.5.1 (Proved by exact exhaustive computation).}
The census in Statement A.1.4 is exhaustive. Moreover the following
finite-range inequality is certified for every real $x$ in its stated
range:
\begin{equation}\label{A:finite-bound}
 |R(x)-C_0x|<\frac{43}{250}\sqrt{x}
 \qquad(10^4\leq x\leq10^9).
\end{equation}
This is a finite theorem, and asserts no bound as $x\to\infty$.
\begin{proof}
The principal census program \texttt{A\_census.cpp} uses independent
$n$-slices and verifies every output in $(s,t,z)$ coordinates. For
$n\geq1$, factor $18n^2-2=2(3n-1)(3n+1)$ with an exact smallest-prime-factor
sieve. In $\mathbb Z[\sqrt3]$, an inert prime of odd exponent rules out a
solution. An inert prime of exponent $2e$ contributes the scalar $p^e$;
the ramified prime $2$ contributes $(1+\sqrt3)^e$; and a split prime of
exponent $e$ contributes the $e+1$ generators
$\pi^k\bar\pi^{e-k}$, $0\leq k\leq e$.
Here $\pi$ is computed as the norm-Euclidean gcd of $p$ and
$r+\sqrt3$, where $r^2\equiv3\pmod p$ is found by the exact
Tonelli--Shanks algorithm. The norm of this gcd is asserted to be $\pm p$.
There is no factor at $3$, since $3n\pm1$ are prime to $3$.
The Euclidean division rounds the two rational coefficients; the strict
decrease of absolute norm is asserted at each division step.

Every product is reduced by exact powers of $2+\sqrt3$, with its real
sign decided by integer comparisons. Since the total norm is positive,
the resulting seeds are precisely
\[
 a+q\sqrt3\in[\sqrt{18n^2-2},\varepsilon_0\sqrt{18n^2-2}),
 \qquad a>2q\geq0.
\]
Unique factorization of ideals and Statement A.1.2 prove that there is
one seed for every unit orbit and no omitted orbit. Equality of distinct
computed seeds is also explicitly ruled out. From each seed, multiply
successively by $\varepsilon_0$ until $4a+17n>X$. These are all solutions
with $a>0,q\geq0$: they correspond to all generators at least
$\sqrt{18n^2-2}$, and the positive coefficient $a$ increases strictly
under this iteration. The $n=0$ family is generated separately from
$1+\sqrt3$ and the strict inequality rejects its first member.
The last possible $n$ is found by exact binary search using
$4\lceil\sqrt{18n^2-2}\rceil+17n\leq X$; at $X=10^9$ it is
$29\,437\,251$. Integer square roots are corrected by exact comparisons,
so floating-point initial guesses do not decide any predicate.

For every emitted point the code checks $F(s,t,z)=-2$,
$t+4z<s$, $4s\leq17z$, the inverse parity, and the original equation
$4(xy+xz+yz)-x^2-y^2-z^2=-1$ in signed 128-bit arithmetic.
Its sorted height list is compared literally with a direct scan of
$s$ at every $z\leq10^6$. A second program,
\texttt{A\_height\_census.cpp}, independently slices by every height
$z\leq10^7$, enumerates ideals of $18z^2-2$, tests the moving region
and the fixed arc $288s^2\leq289(18z^2-2)$, and verifies
\eqref{A:arc-exact} at every height. Its full ordered height list agrees
with the $n$-slice program. This is the required independent recount in
$(s,t,z)$ through $10^7$, and not an extrapolation of the literal scan.

The $10^9$ census contains $35\,363\,970$ solutions at
$23\,243\,639$ occupied heights. The maximum multiplicity is $52$,
first attained at $971\,256\,191$. The raw sorted little-endian
32-bit height receipt has SHA256
\begin{center}\small\ttfamily
29bc345c303f0f1885b13f459669411fc4202\\
a43d3ed95509844f8fac2c65a48
\end{center}
The lossless file \texttt{A\_census\_1e9.d16m8.gz} stores successive
height gaps as 16-bit integers and multiplicities as 8-bit integers;
the maximum gap is $601$. The attached pack/unpack program reconstructs
the raw file and verifies this hash.

To certify \eqref{A:finite-bound}, the Arb computation
\path{A_certify_constant.py}, at 256-bit precision, proves
\[
 \frac{35362047799}{10^{12}}<C_0<
 \frac{35362047800}{10^{12}}.
\]
For both rational endpoints $c/10^{12}$, at every jump height $z$ we
check both its right value and its left limit by the integer inequality
\[
 62500\big(10^{12}R-cz\big)^2
       <1849z\,10^{24}.
\]
The program checks a strict size guard before squaring, which keeps all
products inside signed 128-bit range. It also checks the endpoints
$10^4,10^9$. Convexity in $C_0$ gives the same bound for the enclosed
constant. Between consecutive jumps $R(x)=r$ is constant and
\[
 \frac{d}{dx}\frac{r-C_0x}{\sqrt x}
 =-\frac{r}{2x^{3/2}}-\frac{C_0}{2\sqrt x}<0.
\]
Thus the two one-sided jump values and the external endpoints control
the absolute value throughout every interval. All checks pass, proving
the assertion for real, as well as integer, $x$.
\end{proof}

\paragraph{Statement A.5.2 (Proved finite computations; descriptive evidence with its scope specified).}
On integer $X\leq10^8$ the extremes in Lead 5 are confirmed:
\[
 E(32\,461\,294)=-812.8300691488068897\ldots,
 \qquad E(77\,078\,404)=721.7934249571297570\ldots,
\]
where $E(X)=R(X)-C_0X$. On integer $X\leq10^9$ the extrema are
\[
 \begin{split}
 E(889\,576\,590)&=-2332.8971024648566233\ldots,\\
 E(586\,293\,324)&= 2668.4520483254967593\ldots.
 \end{split}
\]
For $10^4\leq X\leq10^9$ the largest absolute standardized discrepancy
at an integer is $0.1714001928088145597\ldots$, attained at
$X=867\,240$. These decimal evaluations use 80-digit arithmetic after
exactly locating the integer candidates; the finite bound itself is the
separate rational certificate in Statement A.5.1.

The deterministic sample $X=10^4+1000k\leq10^9$, with $999\,991$
observations, has mean $0.00067195$, standard deviation $0.04345364$,
skewness $0.21623$, and excess kurtosis $-0.71874$ for $E(X)/\sqrt X$.
Its median is $-0.00410075$, its $5\%$ and $95\%$ quantiles are
$-0.06359556$ and $0.07133799$, and its lag-one correlation is
$0.99986553$. Thus the sample has strong serial dependence; these
statistics do not constitute evidence of independent normal sampling,
and no limiting law is inferred.
The complete sample and its generating program are attached.

For the continuous mean-square integral, exact constancy of $R(x)$
between successive jumps reduces each interval of length $g$, beginning
at $a$ with $E(a)=e$, to
\[
 \int_a^{a+g}E(x)^2\,dx=e^2g-eC_0g^2+\frac{C_0^2g^3}{3}.
\]
Summation in long-double arithmetic gives the following descriptive
values; they are not interval-certified decimal values:
\[
\begin{array}{c|c}
 X&X^{-2}\int_X^{2X}|E(x)|^2\,dx\\\hline
 10^4&0.0021358981\\
 10^5&0.0057996192\\
 10^6&0.0048290724\\
 10^7&0.0091159753\\
 10^8&0.0027564219\\
 5\cdot10^8&0.0031277302
\end{array}
\]
The asserted finite computations are reproducible from
\path{A_statistics.py} and its output. No global mean-square,
$\Omega$, or $\Omega_\pm$ theorem follows from this finite table.

\paragraph{Statement A.5.3 (Proved weaker mean-square bound).}
For every $\epsilon>0$,
\[
 \int_X^{2X}|R(x)-C_0x|^2\,dx\ll_\epsilon X^{7/3+\epsilon}.
\]
\begin{proof}
Apply Statement A.4.4 with $\epsilon/2$ uniformly for
$x\in[X,2X]$, square it, and integrate over an interval of length $X$.
\end{proof}

\paragraph{Statement A.5.4 (Open: remaining pointwise conjecture and literature scope).}
The exact square-root conjecture for this count is
\[
 \text{for every }\epsilon>0\text{ there are }C_\epsilon,X_\epsilon>0
 \text{ such that }|R(X)-C_0X|\leq C_\epsilon X^{1/2+\epsilon}
 \quad(X\geq X_\epsilon).
\]
The requested mean-square bound is proved in Statement A.5.5.
Statement A.5.8 proves both
$R(X)-C_0X=\Omega_+(X^{1/2})$ and
$R(X)-C_0X=\Omega_-(X^{1/2})$, and hence the unsigned assertion.
The displayed pointwise square-root conjecture remains unproved,
and we have not established an implication from a numbered
published conjecture to it.
The numerical work above settles the finite census and tests the
standardized discrepancy; the infinite oscillation assertions
are settled by the separate analytic proof.

\emph{An exact sufficient arithmetic hypothesis.}
Retain the ideal characters and coefficients $a_\ell(m)$ of
Statement A.4.2, and put
\[
 U_\ell(X)=\sum_{1\leq z\leq X}a_\ell(18z^2-2).
\]
The following uniform estimate would imply the displayed
square-root conjecture:
\begin{equation}\label{A:sufficient-uniform-quadratic}
 \begin{aligned}
 \left|U_\ell(X)-\mathbf1_{\ell=0}\,cX\right|
 &\ll_\eta X^{1/2+\eta}(1+|\ell|)^\eta,\\
 &0\leq|\ell|\leq\lfloor X^{1/2}\rfloor,\qquad \eta>0.
 \end{aligned}
\end{equation}
The implied constant must be independent of both $X$ and $\ell$.
This is a precisely stated sufficient hypothesis, not a conjecture
attributed to the literature. To prove the implication, view the
ideals of norm $18z^2-2$, $z\leq X$, as a positive measure on
$\mathbb R/\mathbb Z$. Its mass is $U_0(X)$ and its $\ell$th
Fourier coefficient is exactly $U_\ell(X)$. Weighted
Erd\H{o}s--Tur\'an, or the Selberg majorants and minorants already
used in Statement A.4.4, gives for $H=\lfloor X^{1/2}\rfloor$
\[
 \left|\sum_{z\leq X}\#\{\mathfrak a:
 N\mathfrak a=18z^2-2,\ \theta(\mathfrak a)\in I\}
       -\lambda_0 U_0(X)\right|
 \ll \frac{U_0(X)}{H+1}
       +\sum_{\ell=1}^{H}\frac{|U_\ell(X)|}{\ell}.
\]
Endpoint conventions change only the already counted Pell
families, hence contribute $O(\log X)$.
Under \eqref{A:sufficient-uniform-quadratic} the right side is
$O_\eta(X^{1/2+3\eta/2}\log(2X))$.
Take $\eta=\epsilon/3$, absorb the logarithm into the remaining
power, and use the exact formula in Statement A.1.3 and
$C_0=\lambda_0c$. This proves the claimed implication. Equivalently,
the nonzero modes in \eqref{A:sufficient-uniform-quadratic} are
\[
 U_\ell(X)=(-1)^\ell
       \sum_{z\leq X}a_\ell(3z-1)a_\ell(3z+1);
\]
thus uniformity in a family of real-quadratic ideal characters
growing as fast as $|\ell|\leq X^{1/2}$ is essential.

The closest explicit shifted-convolution expectation located,
\cite[equation (2) and the paragraph immediately after it]
{NordentoftPetridisRisager2022}, concerns a fixed holomorphic
level-one form and expressly leaves its weight-uniform range
unspecified. It does not state
\eqref{A:sufficient-uniform-quadratic}.
The varying-form Conjectures 1.6 and 1.8 of
\cite{Nelson2022Quadratic} assume non-dihedral holomorphic forms
and irreducible quadratic polynomials, and ask for weaker
cancellation; Remark 1.9 says that variants for dihedral forms
or reducible polynomials still need to be formulated.
Our character family is dihedral and
$18z^2-2=2(3z-1)(3z+1)$ is reducible.
Finally, \cite[Conjecture $C_a(\epsilon)$]{Merikoski2025Primes}
has explicit epsilon dependence on the real-quadratic character
parameter, but its sum runs over the two-variable binary cubic
$c_1y_1^3-c_2y_2^3$. No specialization of that conjecture to
\eqref{A:sufficient-uniform-quadratic} is proved here.

A direct delta-method attempt using the theta-corrected
half-integral Kloosterman-sum Conjecture 1.4 of
Sun \cite{A5Sun2025} was tested on its full analytic weights,
including their inverse-phase amplitudes.  The calculation in
\path{records/B_A5_compensated_weight.tex}, with its exact and
Arb certificates, uses a cutoff supported near the actual quadric
and the dual vectors $(1,0,k)$.  On the fixed modulus block
$2X/3\leq q\leq3X/4$, their normalized weights satisfy
$\|W_k\|_\infty\asymp k^{-1}$ but
$\operatorname{Var}W_k=L+O(k^{-1})$ for an explicit $L>0$.
The variation summed over $k\leq H$ is therefore $\asymp H$,
also for smoothing of width $H^{-1}$.
Consequently, applying the proposed untwisted bound separately
to each frequency by Abel summation and then summing absolutely
has the scale $X^{1/2+\epsilon}H$ on this family.
Balancing this with a boundary cost $X/H$ gives a $3/4$
envelope for that procedure.  This is neither a full $3/4$
theorem nor a proof that Sun's conjecture cannot imply the
desired result by another argument.  It identifies the precise
remaining loss; cancellation between dual frequencies or an
arithmetic transformation retaining their joint phases is still
needed.  The multiplier, bad-prime factors, and complete theta
projection must also be matched before that conjecture is applied.

Completing the normal frequency before summing over moduli gives
a further exact reduction, recorded in
\path{records/A_normal_completion.tex}.
For the same family $(1,0,k)$, a finite quadratic Fourier
transform converts the Sali\'e sums to
$S(2-18h^2,-\overline4;q)$, and Chinese remaindering gives
\[
 S(2-18h^2,-\overline4;q)
       =-S(2-18h^2,-1;2q).
\]
The endpoint phase centers the new index at $h=X$ with width
$X/K$ for an original dyadic block $k\asymp K$.
The bounded-twist Linnik--Selberg conjecture proposed by
Browning--Kumaraswamy--Steiner and stated as
\cite[Conjecture 2]{A5Steiner2019} does apply to each such
completed modulus sum: its fixed parameters can be taken
as $s=24$ and $B=3$. Its resulting bound, summed absolutely
over the new index, is $O_\epsilon(X^{1+\epsilon}/K)$
for this top-face family when $X^\delta\leq K\leq X^{1/2}$,
with any fixed $\delta>0$. Before completion the proposed
Sun estimate gives $O_\epsilon(X^{1/2+\epsilon}K)$.
Taking the better bound still leaves the intermediate
block $K=X^{1/4}$ at scale $X^{3/4+\epsilon}$.
The two exact Fourier identities have been checked in
200 cyclotomic-ring computations, including composite
and prime-power moduli. This supplies a verified application
of a named conjecture to an isolated transformed family;
it supplies no implication to the full sharp square-root
bound. The missing estimate is joint cancellation over
$h,q$, with $2-18h^2$ running over a short quadratic sequence.

The exact scope of the closest sources matters. Theorem 1.5 of
Lekkas--Petridis \cite{LekkasPetridis2026} proves
$\int_X^{2X}|E(x,l)|^2\,dx\ll X^2\log^2X$ for the count integrating
an entire closed geodesic; its Conjecture 1.6 is the square-root
conjecture for that full-geodesic error. Neither statement includes a
fixed sharp foot arc. In \cite[Theorem 3]{Voskou2025Omega}, the
$\Omega_\delta(X^{1/2}(\log\log X)^{1/4-\delta})$ theorem concerns a
full geodesic paired with itself, with the eigenvalue-$1/4$ contribution
removed. Its squared periods supply a positivity absent from the
cross-geodesic count with a source arc. It does not assert
$\Omega_\pm$ for the present count. Statements A.4.4 and A.5.5 prove
the uniform spectral estimates and the requested mean square.
Statement A.5.6 further proves an absolutely convergent
almost-periodic expansion after fixed positive radial smoothing
of the actual sharp arc. It proves both oscillation signs conditional
on a precisely stated nonzero complete-eigenspace coefficient.
Statement A.5.8 proves that condition, indeed for infinitely many
spectral parameters, using the exact arithmetic endpoints,
the outgoing physical kernel, a summable endpoint density, and
the positive shortest-perpendicular contribution.
The square-root pointwise conjecture and its implication from
a numbered known conjecture remain open here.

Two recent results do not remove this gap.
Constantinescu--Nordentoft \cite[Theorems 1.3 and 1.5]{ConstantinescuNordentoft2025}
prove nonvanishing while varying the closed geodesic for a fixed
automorphic form; their general theorem assumes a cusp.
Chatzakos--Dimakis \cite[Theorem 1.2]{ChatzakosDimakis2026} prove
square-root cancellation, up to logarithms, for almost every
point paired with a fixed point on a compact quotient.
Their Section 1.6 explicitly excludes a conclusion for a
preassigned pair. Neither theorem supplies the fixed-cycle,
fixed-arc contraction considered here.


\paragraph{Statement A.5.5 (Proved sharp-arc mean square).}
For every $\epsilon>0$,
\begin{equation}\label{A:sharp-mean-square}
 \int_X^{2X}|R(x)-C_0x|^2\,dx\ll_\epsilon X^{2+\epsilon}.
\end{equation}
\begin{proof}
We extend the uniform transform of Statement A.4.4 to a negative
Mellin line and then apply the geodesic large sieve to its actual
arc coefficients. This checks the additional step needed to use that
large sieve here.

\emph{1. The geometric correction on a negative line.}
In \eqref{A:F-spherical-integral} put
$p=(s+i\kappa)/2$, $B(x)=\cosh(2\delta)+x\sinh(2\delta)$,
and $f(x)=B(x)^{-p}$, where $\delta=\operatorname{arsinh}(1/A)$.
For $\operatorname{Re}s>0$ the identity
\[
 (1-x^2)^{s/2-1}
 =\frac{s+1}{s}(1-x^2)^{s/2}
       -\frac1s\frac d{dx}\{x(1-x^2)^{s/2}\}
\]
and integration by parts give
\begin{equation}\label{A:F-regularized}
 F_{s,\kappa}(A)=
 \frac{\Gamma((s+3)/2)}{\sqrt\pi\Gamma((s+2)/2)}
 \int_{-1}^1(1-x^2)^{s/2}
       \left(f(x)+\frac{x f'(x)}{s+1}\right)dx.
\end{equation}
The right side converges absolutely and is analytic for
$\operatorname{Re}s>-2$, apart from its displayed parameter
singularities. Therefore it continues the left side to the range
$-1<\operatorname{Re}s<0$ used below.

Fix $-1<\sigma<0$, write $s=\sigma+i\tau$,
$U=1+|\tau|$ and $V=U+|\kappa|$. The gamma prefactor is
$O_\sigma(U^{1/2})$, $|s+1|^{-1}\ll_\sigma U^{-1}$,
and $|f'(x)|\ll_\sigma V/A$. Hence
\begin{equation}\label{A:F-negative-large}
 |F_{s,\kappa}(A)-1|
       \ll_\sigma U^{1/2}\left(1+\frac{V}{UA}\right).
\end{equation}
If $A\geq V$, differentiate the integrand in
\eqref{A:F-regularized} twice with respect to $\delta$ on
$0\leq\delta\leq\operatorname{arsinh}(1/A)$.
Its first derivative at zero is odd in $x$, so it integrates to
zero. Its second derivative is $O_\sigma(V^2)$ times the
integrable absolute weight: the possible third power of $p$
from $f'$ is multiplied by $\delta/(s+1)$ and is bounded by
$O_\sigma(V^2)$ because $V\delta\leq1$.
Taylor's integral remainder therefore gives
\begin{equation}\label{A:F-negative-small}
 |F_{s,\kappa}(A)-1|\ll_\sigma U^{1/2}V^2/A^2
 \qquad(A\geq V).
\end{equation}
Using the count $\#\{A_z\leq y\}\ll y$ with its ideal
multiplicities, sum \eqref{A:F-negative-large} over $A_z\leq V$
and \eqref{A:F-negative-small} over $A_z>V$.
The term $V/(UA)$ is summable with total
$O_\sigma(U^{-1/2}V^{1-\sigma})$ because $\sigma<0$;
the other two terms are
$O_\sigma(U^{1/2}V^{1-\sigma})$ because $\sigma>-1$.
Thus precisely the same estimate as \eqref{A:H-uniform} holds:
\begin{equation}\label{A:H-negative}
 |H_\ell(\sigma+i\tau)|
       \ll_\sigma U^{1/2}(U+|\kappa_\ell|)^{1-\sigma}.
\end{equation}
The series defining $H_\ell$ converges normally in this strip,
since its summands at large $A_z$ have two additional inverse
powers of $A_z$.

At $s=0$ the exact hypergeometric identity is
\[
 F_{0,\kappa}(A)=\cos\bigl(\kappa\operatorname{arsinh}(1/A)\bigr).
\]
It follows either by solving the differential equation in the proof
of \eqref{A:F-spherical-integral} at $s=0$, or by its power
series. In particular,
\begin{equation}\label{A:H-at-zero}
 H_0(0)=0,\qquad |H_\ell(0)|\ll1+|\ell|.
\end{equation}
For the second assertion sum
$\min(2,\kappa_\ell^2/A_z^2)$ using the same linear counting
bound. No estimate for the distribution of individual ideal angles
is being assumed in this step.

\emph{2. A uniform spectral expansion with square-root-size remainder.}
Fix a large base height $Y_0$, and take
$T=H=\lfloor\sqrt{Y_0}\rfloor$, fixed as $Y$ varies over
$[Y_0,2Y_0]$. Let $w=w_T$ be either smooth radial cutoff of
Statement A.4.4. Its Mellin transform continues meromorphically by
\[
 \widehat w(s)=-\frac1s\int_0^\infty w'(u)u^s\,du.
\]
It has a simple pole at zero of residue one and satisfies
\eqref{A:w-Mellin-bound} on any fixed nonzero real-part line.
Choose $\sigma=-1/2+\eta$, where $0<\eta<1/2$ is fixed.

The contour estimate \eqref{A:contour-uniform}, with $T=H$,
remains valid on this line. Here are the precise changes to its
proof. The gamma bound \eqref{A:gamma-full-bound} is unchanged,
with $a=1-\eta/2<1$ and $b=3/4-\eta/2$; its numerator gamma
functions still have fixed real parts strictly between $-1/2$
and $-1/4$, and so have no pole on the line. The low-angular and
comparable-angular integrations use only $a<1$ and remain
identical. In the high-angular range, $K\ll H=T$ and the sum is
\[
 K^{1/2-\sigma}\sum_{J\ll K}J^{\sigma+1/2}
       \ll_\eta K,
\]
because $\sigma+1/2=\eta>0$. The exponentially small tails
remain $O(\log(2T))$. These facts give
$O_\eta(T\log(2H))$ for the remaining contour before multiplication
by $Y^\sigma$. All constants are uniform in $Y\in[Y_0,2Y_0]$.
Allowing negative angular indices changes only a factor of two.

We also include angular index zero. Its nonconstant spectral
part obeys the same estimates, using the low-angular range.
The constant eigenfunction contributes
\[
 p'_0p_{0,0}
 \frac{\Gamma((s+1)/2)\Gamma((s-1)/2)}{\Gamma(s/2)^2}.
\]
On the negative contour this is $O_\eta(1)$ in $\tau$, and
its only pole in the traversed strip is at $s=1$.
Its residue in $M_0$ is $4c/3$, where
$c=9L_0^2/(2\pi^2)$. This follows also directly from
\eqref{A:Mellin-geometry} and the already proved
$\sum_{A_z\leq Y}\rho(m_z)=(c/3)Y+O(Y^{2/3})$;
$H_0$ is analytic at one.

There is no contribution at zero from $\widehat w(s)M_\ell(s)$:
the normally convergent gamma series shows $M_\ell(0)=0$
for every $\ell$. For $\ell\ne0$, the factor $1/\Gamma(s/2)$
gives this zero; for $\ell=0$ the nonconstant terms have three
such factors, and the displayed constant term has two.
In the geometric correction the pole of $\widehat w$ does give
$H_\ell(0)$, whose total against coefficients
$|b_\ell|\ll(1+|\ell|)^{-1}$ is $O(H)$ by
\eqref{A:H-at-zero}. The two exceptional double cosets contribute
only $O(\log(2H))$.
Equation \eqref{A:H-negative} bounds the negative-contour
correction by
\[
 O_\eta\left(Y_0^\sigma
       T^{1/2}(T+H)^{1-\sigma}\log(2H)\right)
 =O_\eta\left(Y_0^{1/2+\eta/2}\log Y_0\right).
\]
The remaining $M$ contour costs only
$O_\eta(Y_0^\sigma T\log(2H))$.
The finite-height contour justification from Statement A.4.4
continues without change; the next spectral poles have real part
$-3/2$, outside the traversed strip.

Consequently, for any real trigonometric polynomial
$B(\theta)=\sum_{|\ell|\leq H}b_\ell e^{2\pi i\ell\theta}$
with $|b_\ell|\ll(1+|\ell|)^{-1}$, the weighted count
\[
 \mathcal N_{w,B}(Y)=\sum_z\sum_{N\mathfrak a=m_z}
                  w(A_z/Y)B(\theta(\mathfrak a))
\]
has the expansion, uniformly on $[Y_0,2Y_0]$,
\begin{align}\label{A:sqrt-remainder-expansion}
 \mathcal N_{w,B}(Y)&=\frac{b_0c}{3}\widehat w(1)Y
 \\ &\quad+\frac14\sum_{j\geq1}p'_j
  \left(r_j\widehat w(s_j)Q_{j,H}Y^{s_j}
       +\overline{r_j\widehat w(s_j)Q_{j,H}}Y^{1-s_j}\right)
 \\ &\quad+O_\eta\left(Y_0^{1/2+\eta/2}\log Y_0\right).
 \nonumber
\end{align}
Here $Q_{j,H}=\sum b_\ell D_{s_j}(0)
p_{j,\ell}/D_{s_j}(\kappa_\ell)$ uses the unnormalized source
periods. The conjugation in the second term is valid because
the eigenfunctions are real and $b_{-\ell}=\overline{b_\ell}$.
The coefficients, including $\widehat w(s_j)$ and $Q_{j,H}$,
are fixed throughout this $Y$ interval.

\emph{3. The large sieve and its hypotheses.}
Theorem 3 of Lekkas--Voskou \cite{LekkasVoskou2024}, with weight
$m=0$, states that for a fixed cofinite Fuchsian group and a
fixed closed geodesic, arbitrary complex coefficients $d_j$, and
$\delta$-separated points $y_\nu\in[Y_0,2Y_0]$,
\[
 \sum_\nu\left|\sum_{|t_j|\leq V}d_jy_\nu^{it_j}
                    \int_lu_j\right|^2
 \ll (V+Y_0\delta^{-1})\sum_{|t_j|\leq V}|d_j|^2.
\]
The continuous-spectrum coefficients vanish on a compact quotient.
The same statement is Theorem 1.3 of \cite{LekkasPetridis2026}.
Apply it on the fixed compact torsion-free cover from Statement A.3.1,
with zero coefficients for all eigenfunctions not lifted from
$\Gamma^+$. The covering and period normalization factors are
fixed and were verified above. Apply it separately to the two
target geodesics and use Cauchy--Schwarz for their sum $p_j'$.
There is no requirement that the target equal the source geodesic
or that the coefficients $d_j$ be periods.

For a finite spectral sum, take equally spaced $y_\nu$ with
mesh $h$, choose $\delta=h/2$, multiply the inequality by $h$,
and let $h\to0$. The Riemann sums converge, giving
\begin{equation}\label{A:continuous-large-sieve}
 \int_{Y_0}^{2Y_0}
       \left|\sum_jd_jY^{it_j}p_j'\right|^2dY
       \ll Y_0\sum_j|d_j|^2.
\end{equation}
For the band $J\leq1+t_j<2J$, take
$d_j=r_j\widehat w(s_j)Q_{j,H}$.
By Statement A.4.3, Stirling's formula, and
\eqref{A:w-Mellin-bound},
\[
 \sum_{j\text{ in band}}|d_j|^2
 \ll J^{-1}(1+J/T)^{-2N}\sum_{j\text{ in band}}|Q_{j,H}|^2
 \ll\log^2(2H)(1+J/T)^{-2N}.
\]
The factor $Y^{1/2}$ in \eqref{A:sqrt-remainder-expansion}
therefore gives $L^2([Y_0,2Y_0])$ norm at most
$O(Y_0\log(2H)(1+J/T)^{-N})$ for this band.
Minkowski's inequality over the dyadic bands, and then passage
to the absolutely convergent spectral series, prove
\begin{equation}\label{A:smoothed-mean-square}
 \int_{Y_0}^{2Y_0}
 \left|\mathcal N_{w,B}(Y)-\frac{b_0c}{3}\widehat w(1)Y\right|^2dY
 \ll_\eta Y_0^2\log^4Y_0+Y_0^{2+\eta}\log^2Y_0.
\end{equation}

\emph{4. Sharp arc and sharp height.}
For the fixed arc $I=[0,\lambda_0]$ take its Selberg majorant
and minorant trigonometric polynomials $B_H^+,B_H^-$.
The exact elementary approximation used here is
\cite[equation (4.11)]{LesterMatomakiRadziwill2015}:
for every interval in $\mathbb R/\mathbb Z$ and integer $H\geq1$,
\[
 B_H^-\leq\mathbf1_I\leq B_H^+,\qquad
 b_0^\pm=|I|\pm\frac1{H+1},\qquad
 |b_\ell^\pm|\ll(1+|\ell|)^{-1}\quad(0<|\ell|\leq H).
\]
Endpoint half-value conventions can change the sum only by the
already proved $O(\log Y_0)$ Pell contribution.
These polynomials are real, have degree at most $H$, and are
fixed throughout the dyadic $Y$ interval. In particular $B_H^+\geq0$.
Use the radial cutoffs $w^\pm$ of Statement A.4.4 with the same
fixed $T=\lfloor\sqrt{Y_0}\rfloor$. Then
\[
 \mathcal N_{w^-,B_H^-}(Y)
 \leq\mathcal A(Y)
 \leq\mathcal N_{w^+,B_H^+}(Y)+O(\log Y_0),
\]
where $\mathcal A(Y)$ is the exact sharp ideal-arc count with
$A_z\leq Y$. The lower inequality is valid even where
$B_H^-$ is negative: if $w^-\ne0$ the radial sharp indicator
is one, and multiplication by $0\leq w^-\leq1$ still leaves
$w^-B_H^-\leq\mathbf1_I$. Outside that radial range it is zero.
The upper inequality follows from nonnegativity of $B_H^+$.

Both main terms differ from $\lambda_0cY/3$ by
$O(Y_0/H+Y_0/T)=O(\sqrt{Y_0})$.
Thus \eqref{A:smoothed-mean-square} for the two fixed pairs
$(w^-,B_H^-)$ and $(w^+,B_H^+)$ proves
\[
 \int_{Y_0}^{2Y_0}|\mathcal A(Y)-\lambda_0cY/3|^2dY
       \ll_\epsilon Y_0^{2+\epsilon},
\]
after choosing $\eta$ smaller than $\epsilon$ and absorbing
logarithms. Finally $Y=(9x^2-1)^{1/2}$ has derivative bounded
above and below by positive constants for $x\geq1$;
two dyadic $Y$ intervals cover the image of $[X,2X]$.
The difference between $\lambda_0cY/3$ and $C_0x$ is
$O(x^{-1})$, and Statement A.1.3 adds $O(\log(2x))$.
These terms have mean square $O(X\log^2X)$, proving
\eqref{A:sharp-mean-square}.
\end{proof}


\paragraph{Statement A.5.6 (Proved smoothed expansion and conditional two-sided oscillation).}
Let
\[
 \mathcal A(Y)=\sum_{A_z\leq Y}
       \#\{\mathfrak a:N\mathfrak a=m_z,\ \theta(\mathfrak a)\in I\},
 \qquad e(Y)=\frac{\mathcal A(Y)-\lambda_0cY/3}{\sqrt Y}.
\]
Fix a smooth nonincreasing function $w:[0,\infty)\to[0,1]$,
equal to $1$ near zero and to $0$ for sufficiently large argument.
Its derivative is supported in a compact subinterval of $(0,\infty)$.
Then, as $Y\to\infty$,
\begin{align}\label{A:fixed-smoothing-AP}
 \int_0^\infty e(Yv)\sqrt v\,(-w'(v))\,dv
 &=\frac14\sum_{j\geq1}p'_j
 \left(r_j\widehat w(s_j)Q_{j,I}Y^{it_j}
 +\overline{r_j\widehat w(s_j)Q_{j,I}}Y^{-it_j}\right)
 \nonumber\\
 &\quad+O_w(Y^{-1/25}).
\end{align}
The series converges absolutely and uniformly in $\log Y$.
It is a real uniformly almost-periodic function with mean zero.
The following is consequently a sufficient condition for the
requested two-sided oscillation:
\begin{equation}\label{A:exact-nonvanishing-condition}
 \text{there is a }t>0\text{ such that }\quad
 \mathcal C(t):=\sum_{j:t_j=t}p'_jQ_{j,I}\ne0.
\end{equation}
If \eqref{A:exact-nonvanishing-condition} holds, then
$R(X)-C_0X=\Omega_+(X^{1/2})$ and
$R(X)-C_0X=\Omega_-(X^{1/2})$.
The sum in \eqref{A:exact-nonvanishing-condition} is over the
complete eigenspace and is independent of its real orthonormal
basis. Nonvanishing is not asserted in this statement.

\begin{proof}
We give the additional estimates needed to pass to the sharp arc.
They also explain why the $O(H)$ zero-residue bound used in
Statement A.5.5 cannot simply be retained in this argument.

\emph{1. Uniform angular discrepancy.}
The proof of Statement A.4.4 gives the stronger intermediate fact
\begin{equation}\label{A:all-interval-discrepancy}
 D(V):=\sup_{J\subset\mathbb R/\mathbb Z}
 \left|\sum_{A_z\leq V}\sum_{N\mathfrak a=m_z}
       \bigl(\mathbf1_J(\theta(\mathfrak a))-|J|\bigr)\right|
 \ll_\epsilon V^{2/3+\epsilon}.
\end{equation}
Indeed its phase-aligned bound for
$\sum_{\ell\leq H}|W_\ell|/\ell$ is independent of the interval;
the weighted Erd\H{o}s--Tur\'an step and the two radial
brackets therefore apply to every $J$ with the same constants.
One can use half-open intervals, or control a point mass by
the same positive Fej\'er majorant. This avoids assuming any
uniform Pell estimate for a varying interval endpoint.
For every periodic absolutely continuous function $f$ of mean zero,
Stieltjes integration by parts consequently gives
\begin{equation}\label{A:BV-discrepancy}
 \left|\sum_{A_z\leq V}\sum_{\mathfrak a}f(\theta(\mathfrak a))\right|
 \leq D(V)\operatorname{Var}(f).
\end{equation}
The assertion follows first for step functions, by expressing
them through interval indicators, and then by approximation.

\emph{2. Canonical angular polynomials have bounded variation.}
We use a specified choice of the Selberg polynomials.
The sawtooth approximation in
\cite[Lemma 2.2]{WangZhang2023}, explicitly attributed there to
Vaaler's Theorem 18, states the following. For an integer $H\geq1$,
put
\[
 W(t)=\pi t(1-|t|)\cot(\pi t)+|t|\quad(0<|t|<1),
 \qquad W(0)=1,
\]
extend $W$ by zero outside $[-1,1]$, and set
\[
 \psi_H(x)=-\sum_{1\leq|\ell|\leq H}
       \frac{W(\ell/(H+1))}{2\pi i\ell}e^{2\pi i\ell x},
 \quad
 F_H(x)=\sum_{|\ell|\leq H}
       \left(1-\frac{|\ell|}{H+1}\right)e^{2\pi i\ell x}.
\]
Then $F_H\geq0$ and
$|\psi_H-\psi|\leq F_H/(2(H+1))$.
These hypotheses are purely trigonometric and impose no
condition on our lattice. For an interval $I=[a,b]$ take
\[
 B_H^\pm(x)=b-a+\psi_H(a-x)+\psi_H(x-b)
 \ \pm\frac{F_H(a-x)+F_H(x-b)}{2(H+1)}.
\]
With the usual half-value convention at the two endpoints,
these satisfy the pointwise brackets of Statement A.5.5.
The endpoint convention changes our actual count by
$O(\log(2Y))$, already proved in Statement A.1.3.

In addition to the previously used Fourier bounds, these
particular polynomials satisfy
\begin{equation}\label{A:Selberg-BV}
 \|B_H^\pm\|_\infty+\operatorname{Var}(B_H^\pm)\ll1,
 \qquad \|(B_H^\pm)^{(r)}\|_1\ll_r H^{r-1}\quad(r\geq1).
\end{equation}
Here is a direct verification of the extra assertion.
The extension of $W$ is $C^1$, with integrable piecewise second
derivative: its value and first derivative vanish at $\pm1$.
Twice summing its Fourier polynomial by parts gives
\[
 \left|\sum_{\ell\in\mathbb Z}W(\ell/(H+1))e^{2\pi i\ell x}\right|
 \ll \min\{H,(H\|x\|^2)^{-1}\}.
\]
Its $L^1$ norm is therefore bounded independently of $H$.
Since $\psi_H'$ equals one minus this polynomial, it has
bounded $L^1$ norm. The formula
$F_H(x)=\sin^2(\pi(H+1)x)/((H+1)\sin^2(\pi x))$
shows $\|F_H'\|_1\ll H$: split at $\|x\|=1/H$ and
differentiate, using $O(H^3|x|)$ in the first range and
$O(|x|^{-2}+H^{-1}|x|^{-3})$ in the second.
This proves the variation bound; bounded sup norm follows from
the sawtooth approximation and $F_H\leq H+1$.
Finally, the elementary $L^1$ Bernstein bound for band-limited
polynomials gives the derivative estimates in
\eqref{A:Selberg-BV}. It can be proved here without an
approximation theorem: choose a fixed smooth frequency cutoff
equal to one on $[-1,1]$ and supported in $[-2,2]$.
Its periodic kernel, differentiated $r-1$ times, has $L^1$
norm $O_r(H^{r-1})$, by periodizing its Schwartz inverse Fourier
transform. Convolve this kernel with $(B_H^\pm)'$.

\emph{3. Improvement of the zero residue.}
For either polynomial $B=B_H^\pm$, the exact identity at $s=0$
in Statement A.5.5 gives
\[
 H_B(0):=\sum_{|\ell|\leq H}b_\ell H_\ell(0)
   =\sum_z\sum_{\mathfrak a}f_{A_z}(\theta(\mathfrak a)),
\]
where, putting $\delta(A)=\operatorname{arsinh}(1/A)/L_0$,
\[
 f_A(x)=\tfrac12\{B(x+\delta(A))+B(x-\delta(A))-2B(x)\}.
\]
This series converges absolutely for each fixed $H$,
since its large-$A$ summand is $O_H(A^{-2})$ and
$S_A(V)\ll V$. Every $f_A$ has mean zero.
On a dyadic block $A\asymp Z$, \eqref{A:Selberg-BV} gives
\[
 \operatorname{Var}(f_A)\ll
      \min\{1,H^2/Z^2\},\qquad
 \operatorname{Var}(\partial_A f_A)
      \ll\min\{H/Z^2,H^2/Z^3\}.
\]
The second estimate follows from
$|\delta'(A)|\ll A^{-2}$ and the first or second difference
of $B''$; the first follows from the second difference of $B'$.
Partial summation in $A$, followed by
\eqref{A:BV-discrepancy}, bounds that dyadic block by
\[
 \begin{cases}
 O_\epsilon(HZ^{\vartheta-1}),&Z\leq H,\\
 O_\epsilon(H^2Z^{\vartheta-2}),&Z\geq H,
 \end{cases}
 \qquad \vartheta=2/3+\epsilon<1.
\]
For $A\leq V$ use instead $\|B\|_\infty\ll1$ and
$S_A(V)\ll V$. Summing the remaining dyadic blocks gives
\[
 |H_B(0)|\ll_\epsilon
 V+HV^{\vartheta-1}+H^\vartheta.
\]
Take $V=H^{1/(2-\vartheta)}$. Since $\epsilon$ can be
arbitrarily small, this proves
\begin{equation}\label{A:zero-residue-improved}
 H_B(0)\ll_\epsilon H^{3/4+\epsilon}.
\end{equation}

\emph{4. Fixed radial smoothing and passage to the sharp arc.}
Use the negative-line transform of Statement A.5.5 with the fixed
function $w$, $\sigma=-49/100$, and $H=\lfloor Y^{3/5}\rfloor$.
The parameter $T$ in the transform bounds is now a fixed
constant depending on $w$. The contour bound remains
$O_w((T+H)\log(2H))$ for $-1/2<\sigma<0$.
To verify the only extra range not used in Statement A.5.5,
take $K>T$ in its high-angular calculation. The dyadic sum is
\[
 K^{1/2-\sigma}\sum_{J\ll K}J^\sigma\min(J,T)^{1/2}
 \ll_\sigma K^{1/2-\sigma}T^{\sigma+1/2}\leq C_\sigma K.
\]
The low and comparable ranges use only $a<1$, exactly as before.
The spherical-correction inversion, the new zero-residue bound,
and the exceptional intersection terms therefore leave an error
\[
 O_w\bigl(Y^\sigma H\log(2H)
          +Y^\sigma H^{1-\sigma}\log(2H)
          +H^{3/4+1/100}+\log(2Y)\bigr)
   =O_w(Y^{23/50}).
\]
The zero Fourier mode has the same main term and pole at $s=1$
as in Statement A.5.5; all $M_\ell(0)$ vanish.

We also need a bound uniform in $H$ for the individual
$Q_{j,H}^\pm$. The quantitative assertion of
\cite[Theorem 7.1 and Section 7.1]{AnantharamanZelditch2012}
gives a polynomial bound in $t_j$ for the H\"older norm of
the boundary primitive. Its displayed constant is
$O((1+t_j)^{3/2})$ times
$\|u_j\|_\infty+\|\nabla u_j\|_\infty$.
On the fixed compact cover, the elliptic estimate and Sobolev
embedding give the deliberately coarse bound
$\|u_j\|_{C^1}\ll(1+t_j)^4$.
The fixed coordinate changes and multiplication by
$(2e^v)^{-s_j}$ in Statement A.4.6 cost at most a further fixed
power of $1+t_j$. Thus, for some absolute integer $M$ depending
only on this fixed geometric setup, the periodic primitive
$F_j$ of $\beta_j$ has
\[
 \|F_j\|_{C^{1/2}}\ll(1+t_j)^M.
\]
An affine mean is understood in this assertion, as in
Statement A.4.6.

Identify $u=L_0\theta$, absorbing only fixed coordinate factors.
The middle polynomial in $B_H^\pm$ is the convolution of the
interval indicator with the kernel having coefficients
$W(\ell/(H+1))$. Its bounded $L^1$ norm and the displayed
decay bound imply approximation of a $1/2$-H\"older primitive
with error $O(H^{-1/2})$. For the two Fej\'er terms,
integration by parts and the preceding bounds for $F_H'$ give
\[
 \frac1{H+1}\left|\int F_H(x-a)\,\beta_j(dx)\right|
 \ll (1+t_j)^M H^{-1/2};
\]
subtract $F_j(a)$ inside the derivative integral and use
$|F_j(x)-F_j(a)|\ll(1+t_j)^M\|x-a\|^{1/2}$.
Consequently
\begin{equation}\label{A:Selberg-Q-limit}
 |Q_{j,H}^\pm-Q_{j,I}|
       \ll(1+t_j)^M H^{-1/2},
 \qquad |Q_{j,H}^\pm|+|Q_{j,I}|\ll(1+t_j)^M,
\end{equation}
after increasing $M$ if necessary.
The rapid vertical decay of $\widehat w$, the polynomial gamma
bound on $r_j$, the broad-band period bound, and
$\#\{t_j\leq V\}\ll V^2$ make the residue series absolutely
convergent. Equation \eqref{A:Selberg-Q-limit} permits passage
to its sharp-arc limit with error $O_w(Y^{1/2}H^{-1/2})$.

Bracket the actual radially smoothed arc count by $B_H^-$ and
$B_H^+$. Their main terms differ from
$\lambda_0cY\widehat w(1)/3$ by $O_w(Y/H)$.
Both spectral terms converge to the same series, and all the
other errors just proved are $O_w(Y^{23/50})$.
This proves the corresponding unnormalized expansion.
Partial summation gives
\[
 \sum_z\sum_{\mathfrak a}w(A_z/Y)\mathbf1_I(\theta(\mathfrak a))
       -\frac{\lambda_0c}{3}Y\widehat w(1)
 =\sqrt Y\int_0^\infty e(Yv)\sqrt v\,(-w'(v))\,dv,
\]
and proves \eqref{A:fixed-smoothing-AP}.

\emph{5. The conditional two-sided conclusion.}
Group the absolutely convergent series by distinct $t$.
Its coefficient at frequency $t$ is
$r(t)\widehat w(1/2+it)\mathcal C(t)/4$.
Here $r(t)\ne0$ because the gamma function has no zeros.
If \eqref{A:exact-nonvanishing-condition} holds, choose $w$
with $\widehat w(1/2+it)\ne0$: smooth monotone approximations
to $\mathbf1_{[0,1]}$ have Mellin transforms tending to
$(1/2+it)^{-1}\ne0$.
The resulting real series $P(\log Y)$ is nonzero.
Termwise averaging, justified by absolute uniform convergence,
gives mean zero and strictly positive mean square.
Since it is bounded, these two facts imply
$\limsup P>0$ and $\liminf P<0$.
For example, if $\limsup P\leq0$, then for every $\varepsilon>0$
one has eventually $P\leq\varepsilon$ and $P\geq-C$.
The inequality
$P^2\leq(\varepsilon-C)P+\varepsilon C$ contradicts the
positive mean square after averaging and letting
$\varepsilon\to0$. The other sign is identical.

The kernel $\sqrt v\,(-w'(v))$ is nonnegative, has finite
positive integral, and is supported in a fixed compact
subinterval of $(0,\infty)$. Positive and negative values
bounded away from zero of its convolution with $e(Y)$
therefore force positive and negative values bounded away
from zero of $e$ at arbitrarily large arguments.
Finally $Y=\sqrt{9X^2-1}=3X+O(X^{-1})$ and
$R(X)=\mathcal A(Y)+O(\log(2X))$ transfer both assertions
to $R(X)-C_0X$. The eigenspace sum is basis-independent because
it is the contraction of the two linear period functionals
with the orthogonal projector onto that eigenspace.
\end{proof}
